\RequirePackage{setspace}
\documentclass{jgea}

\volume{} 		% Please input the volume of the journal.
\pubyear{}		% Please input the publication year.
\issueno{} 		% Please input the issue number.
\pagerange{}{} 	% Please input the first and last page numbers.
\setcounter{page}{1}

%	-------------------------------------------------------------------
%	Files included by this document live in ./jgea_new/ (figures,
%	appendices, bibliography, generated data). The figure files in
%	./jgea_new/figures/ are written by ./jgea_new/make_figs.py from the
%	results in ../forensics/figures/ -- rerun that script after new
%	model runs.
%	-------------------------------------------------------------------

\usepackage{rotating}
\usepackage{xspace}
\usepackage{tikzscale}
\usepackage{tikz}
\usepackage{fancyvrb}
\usepackage{subcaption}
\usepackage{pgfplotstable}
\usepackage{longtable}
\usepackage{pgfplots}
\usetikzlibrary{arrows.meta}
\usepackage{amsmath}
\usepackage{hyperref}
\usepackage{helvet}
\usepackage[eulergreek]{sansmath}

\pgfplotsset{compat=1.3}

\fvset{formatcom=\color{black},fontsize=\scriptsize,frame=single,rulecolor=\color{black},baselinestretch=1,framesep=5mm}
\renewcommand{\FancyVerbFormatLine}[1]{\textcolor{blue}{#1}}

\newcommand{\GDX}{\textsc{gdx}\xspace}
\newcommand{\HAR}{\textsc{har}\xspace}
\newcommand{\RHS}{\textsc{rhs}\xspace}
\newcommand{\LHS}{\textsc{lhs}\xspace}
\newcommand{\CES}{\textsc{ces}\xspace}
\newcommand{\CET}{\textsc{cet}\xspace}
\newcommand{\CD}{\textsc{cd}\xspace}
\newcommand{\CDE}{\textsc{cde}\xspace}
\newcommand{\LES}{\textsc{les}\xspace}
\newcommand{\CGE}{\textsc{cge}\xspace}
\newcommand{\SOE}{\textsc{soe}\xspace}
\newcommand{\LOE}{\textsc{loe}\xspace}
\newcommand{\SATE}{\textsc{sate}\xspace}
\newcommand{\AUES}{\textsc{aues}\xspace}
\newcommand{\GMR}{\textsc{gmr}\xspace}
\newcommand{\GMRLS}{\textsc{gmr\_ls}\xspace}
\newcommand{\MCP}{\textsc{mcp}\xspace}
\newcommand{\MGE}{\textsc{mge}\xspace}
\newcommand{\IOT}{\textsc{iot}\xspace}
\newcommand{\SAM}{\textsc{sam}\xspace}
\newcommand{\EOS}{\textsc{eos}\xspace}
\newcommand{\EOT}{\textsc{eot}\xspace}
\newcommand{\USA}{\textsc{usa}\xspace}
\newcommand{\CHN}{\textsc{chn}\xspace}
\newcommand{\ZAF}{\textsc{zaf}\xspace}
\newcommand{\ARG}{\textsc{arg}\xspace}
\newcommand{\HKG}{\textsc{hkg}\xspace}
\newcommand{\GAMS}{\textsc{gams}\xspace}
\newcommand{\MPSGE}{\textsc{mpsge}\xspace}
\newcommand{\GTAP}{\textsc{gtap}\xspace}
\newcommand{\GTAPINGAMS}{\textsc{gtapingams}\xspace}
\newcommand{\GTAPtwelve}{\textsc{gtap12ingams}\xspace}
\newcommand{\GTAPnine}{\textsc{gtap9ingams}\xspace}
\newcommand{\GEMPACK}{\textsc{gempack}\xspace}
%	-------------------------------------------------------------------

\begin{document}
\title{A Small Armington Trade Economy: \\ Terms of Trade in Single-Country Models, Evaluated with \GTAPtwelve} \author{ Florian Landis\footnote{ZHAW Zurich University of Applied Sciences, Winterthur, Switzerland (email: ladi@zhaw.ch).} }

\maketitle

\begin{abstract}

Single-country computable general equilibrium models usually rest on the small open economy (\SOE) assumption that world prices of both imports and exports are fixed.
Such models cannot represent the terms of trade effects of trade policy, which matter for welfare.
I propose a simple modification, the \emph{small Armington trade economy} (\SATE): import prices remain fixed, but the rest of the world demands the economy's exports with a constant price elasticity that equals the Armington elasticity of substitution between import sources.
The specification follows from the limit of an Armington world economy in which the country becomes small, it requires no data on the rest of the world beyond the country's own exports, and it can be implemented in a few lines of code.
Using an update of the \GTAPINGAMS model of \citet{LanzRutherford2016} to the \GTAP 12 Data Base, I compare the \SATE and the \SOE closures with the global multiregional model for unilateral, proportional changes of trade taxes in the United States, China, South Africa, and Argentina.
The \SATE closure tracks the welfare effects of the global model closely and gets much closer to the welfare-maximizing tariff level than the \SOE closure.
Its deviations from the global model are smaller than the effect of halving the Armington elasticities in the global model.
%Sectoral results are reproduced almost exactly for the two small economies, but less well for the two large ones.
%The paper also documents \GTAPtwelve, the update of \GTAPnine to the \GTAP 12 Data Base.

\end{abstract}

\begin{jelcodes}
C68, F13, F14, F17
\end{jelcodes}

\begin{keywords}
Small open economy; Terms of trade; Armington elasticity; Single-country CGE models; Trade policy; \GTAP.
\end{keywords}

%	===================================================================
\section{Introduction}\label{sec:intro}
%	===================================================================

Restrictions in data availability or model dimensionality imply that many computable general equilibrium (\CGE) models are written as single-country models.
More often than not, the small open economy (\SOE) assumption is chosen for formulating such models: global market prices are taken as fixed.
Without careful calibration, fixed global market prices can make exported quantities very sensitive to domestic price changes.
The classical approach to this problem was given by \citet{DeMeloRobinson1989}.
It works around overly sensitive export quantities by calibrating finite price elasticities for both imports and exports:
The provision of domestically produced goods to local and international markets is subject to imperfect transformation (described by a finite elasticity of transformation, \EOT), just as intermediate and final demand is subject to imperfect substitution between imported and domestically supplied varieties (described by a finite elasticity of substitution, \EOS).
This yields reasonable quantity responses to price shocks, but \citet{LanzRutherford2016} and \citet{Schreiber2024}, as well as the sources cited therein, show that \SOE models calibrated in this way cannot keep track of the welfare effects of trade policy in a way that is consistent with global models, because they miss terms of trade effects.

One known remedy is to switch to a large open economy (\LOE) model.
Developed by James Markusen \citep{Markusen}\footnote{%
  The formulation of the \LOE is attributed to James Markusen and has been around since the 1990s \citep{Schreiber2024}.
  Its availability online is not reliable everywhere.
  I find \url{https://www.mpsge.org/tutorial.pdf} a good source.
  (The formulation of the \LOE model is given in example ``Model 48'' on pages 31 ff.)
}, this framework has found application in models such as the MIT \textsc{usrep} model \citep{Yuan2019}.
But the correct calibration of the price elasticities of the world's demand for the country's exports and of the world's supply of its imports is not self-evident, and the only systematic attempt at this I am aware of is the recent paper by \citet{Schreiber2024}.
They use the \GTAPINGAMS model for in-model estimation of these elasticities and calibrate an \LOE version of the US economy, which they show to be very successful in reproducing the results of the global \GTAPINGAMS model.

In this paper, I go a similar yet simpler route.
I observe that the missing terms of trade effects in \SOE models are a direct consequence of the assumption that both import and export prices are fixed at world market levels: the small economy cannot affect them.
But this is not the situation of small countries in a global \CGE model such as the \GTAP model \citep{Corong2017}.
The Armington assumption implies that country A's output is traded off against imports from elsewhere  at a constant elasticity of substitution between import sources.
As the size of A's economy relative to the world economy goes to zero, the sectoral activity levels and the overall import demand of the other countries become independent of the prices in A.
But the market power of A's exports \emph{within} the other countries' imports does not vanish with A's relative size.
\citet{CaliendoFeenstra2024} show formally that, in the limit of an Armington economy whose size relative to its trading partners goes to zero, the world's demand for its exports has a constant price elasticity equal to the elasticity of substitution between import varieties.

I therefore propose an improved \SOE closure, the \emph{small Armington trade economy} (\SATE), in which the world's demand for the country's exports of commodity $i$ has a constant price elasticity equal to the Armington elasticity $\sigma^m_i$, while import prices remain fixed at exogenous world price levels.
The \SATE closure requires no information about the rest of the world beyond the country's benchmark exports and a set of Armington elasticities.
It is therefore particularly attractive for models based on single-country input--output tables (\IOT{}s) or social accounting matrices (\SAM{}s), rather than on global databases such as \GTAP.

The \GTAPINGAMS framework is well suited to evaluate the \SATE closure, because the same benchmark data, technologies, and preferences can be solved as a global multiregional (\GMR) model and as a small open economy.
I update the \GTAPINGAMS code of \citet{LanzRutherford2016} from \GTAP 9 to the \GTAP 12 Data Base \citep{Aguiar2025}, add two implementations of the \SATE closure, and repeat their illustrative experiment: unilateral, proportional changes of a country's import tariffs and export taxes.
I report results for two large economies, the United States and China, and for the two smallest G20 economies in the dataset, South Africa and Argentina.

The results show that the \SATE closure tracks the welfare effects of the \GMR model closely and identifies the welfare-maximizing level of trade taxes much better than the \SOE closure, which recommends free trade in almost all cases.
The careful calibration of \citet{Schreiber2024} can be expected to be more accurate still, but the deviations of the \SATE results from the \GMR results do not exceed the sensitivity of the \GMR results to uncertain Armington elasticities, which I assess by halving these elasticities.
The straightforward calibration of the \SATE closure makes it an attractive choice for single-country \CGE models.
For sectoral activity levels, the \SATE closure reproduces the global model almost exactly for the two small economies, but tends to overstate output responses for the two large ones.

In the following, I describe the model and the modifications relative to \citet{LanzRutherford2016} (Section \ref{sec:model}), compare the \SATE and \SOE closures with the global model (Section \ref{sec:comparison}) and conclude (Section \ref{sec:concl}).
% Appendix \ref{app:model} documents the model in full, largely reproducing the model documentation of \citet{LanzRutherford2016}.

%	===================================================================
\section{Model Description}\label{sec:model}
%	===================================================================

The model is the static \GTAPINGAMS model of \citet{LanzRutherford2016}, implemented in \GAMS/\MPSGE, updated to the \GTAP 12 Data Base, and extended by the \SATE closure.
This section gives an overview of the main features and emphasizes what has changed.
A complete documentation is given in Appendix \ref{app:model}.

\subsection{Overview of the \GTAPINGAMS model}

Each region hosts a representative agent who owns the primary factors (labor, capital, land, and natural resources), collects all tax revenues, and finances fixed levels of public consumption and investment as well as the benchmark current account balance, which is held fixed.
Production in each sector combines intermediate inputs and a \CES value-added aggregate of primary factors in a Leontief technology; land and natural resources are sluggish and allocated to sectors through a constant elasticity of transformation (\CET) function.
Each intermediate input and each good in final demand is an Armington composite of domestic and imported varieties, with elasticity of substitution $\sigma^{DM}_i$, and the imported variety is itself a \CES composite of imports from all source regions with elasticity $\sigma^m_i$ (\texttt{esubm}).
Bilateral trade flows carry international transport margins, export taxes or subsidies of the exporter, and import tariffs of the importer.
Output of a sector can be sold on the domestic market and exported; a \CET function with elasticity $\eta^{DX}_i$ governs the transformation between the two.
Final demand of the representative agent can be represented by Cobb-Douglas (\CD), linear expenditure (\LES), or constant difference of elasticities (\CDE) preferences.
All simulations reported here use \CD preferences (\citealt{LanzRutherford2016} found the demand system to matter little for the experiment considered here).

The model can be solved as a global multiregional model (\GMR) or as a small open economy model (\SOE), in which only one region (or a subset of regions) is modelled explicitly and the remaining regions form a rest of the world (\textsc{row}) with fixed prices.
In the \SOE closure of \citet{LanzRutherford2016}, exports of the modelled region are converted into foreign exchange at fixed world prices, imports are available in any quantity at fixed world prices, and a real exchange rate, $p^{FX}$, adjusts so as to keep the current account at its benchmark level.
Both closures are calibrated to the same benchmark data and use the same elasticities.
In the \GMR closure, a region's policy affects the prices of its exports in all importing regions and the prices of its imports from all exporting regions; in the \SOE closure, these terms of trade effects are absent by construction.

\subsection{The small Armington trade economy}\label{sec:sate}

Consider a small region $r$ and its exports of good $i$ to an importing region $k$.
With \CES aggregation of imports from all sources at elasticity $\sigma^m_i$, region $k$'s demand for the variety from $r$ is
\begin{equation}\label{bilateral}
x_{irk} = M_{ik}\,\alpha_{irk}\left(\frac{P^M_{ik}}{p_{irk}}\right)^{\sigma^m_i},
\end{equation}
where $M_{ik}$ and $P^M_{ik}$ are the quantity and price of $k$'s import aggregate and $p_{irk}$ is the price of $r$'s variety in $k$.
The price elasticity of this demand is
\begin{equation}\label{elasticity}
-\frac{\partial \ln x_{irk}}{\partial \ln p_{irk}} = \sigma^m_i\,(1-s_{irk}) + s_{irk}\,\epsilon^M_{ik},
\end{equation}
where $s_{irk}$ is the share of $r$ in $k$'s imports of $i$ and $\epsilon^M_{ik}$ is the price elasticity of $k$'s aggregate import demand.
See, e.g., \citet{AtkesonBurstein2008} and their equation (19) for a formal description of the mathematical framework that yields \eqref{elasticity}.
As $r$ becomes small, $s_{irk}\to 0$ for all importers, and the elasticity of the world's demand for $r$'s exports approaches $\sigma^m_i$.
\citet{CaliendoFeenstra2024} derive this limit rigorously for the Armington model (and analogous limits for Eaton--Kortum, Krugman, and Melitz models).
The demand of the rest of the world for the exports of the small economy is therefore not perfectly elastic, while the small economy's imports do have fixed world prices in the limit.
The \SATE closure replaces the perfectly elastic export demand of the \SOE closure by
\begin{equation}\label{sate}
E_{ir} = \left(\frac{p^{FX}}{p^{PX}_{ir}}\right)^{\sigma^m_i},
\end{equation}
where $E_{ir}$ is an index of export quantities (unity in the benchmark) and $p^{PX}_{ir}$ is the price of $r$'s exports of $i$.
Import prices remain fixed in terms of foreign exchange.
% Figure \ref{fig:exportdemand} illustrates the difference between the two closures.

% \begin{figure}[t!]
% \centering
% \begin{tikzpicture}
% \begin{axis}[
%   width=0.6\textwidth, height=0.3\textheight,
%   xlabel={Export quantity (index)}, ylabel={Export price relative to $p^{FX}$},
%   xmin=0, xmax=2.5, ymin=0.6, ymax=1.5, grid=major, grid style={dashed, gray!30},
%   legend style={draw=none, at={(0.98,0.98)}, anchor=north east}, legend cell align=left]
% \addplot[color=blue, dotted, thick, domain=0:2.5] {1};
% \addlegendentry{\SOE: perfectly elastic}
% \addplot[color=red, thick, domain=0.25:2.5, samples=80] {x^(-1/4)};
% \addlegendentry{\SATE: elasticity $\sigma^m_i=4$}
% \addplot[color=red, dashed, domain=0.25:2.5, samples=80] {x^(-1/8)};
% \addlegendentry{\SATE: elasticity $\sigma^m_i=8$}
% \node[circle, fill=black, scale=0.4] at (axis cs:1,1) {};
% \end{axis}
% \end{tikzpicture}
% \caption{World demand for the exports of a small economy in the \SOE and \SATE closures.}\label{fig:exportdemand}
% \end{figure}

Two caveats apply.
First, in the \GTAP model the elasticity in \eqref{elasticity} refers to the importer's price, which includes transport margins and the importer's tariffs; the elasticity with respect to the exporter's f.o.b.\ price is smaller in proportion to the f.o.b.\ share of the delivered price.
Second, the limit argument does not hold for the largest economies, whose shares $s_{irk}$ are far from zero in some markets; for these, \eqref{sate} is an approximation whose quality is an empirical question.
% Both caveats suggest that $\sigma^m_i$ is an upper bound for the true export demand elasticity.

I implement the \SATE closure in two ways (Appendix \ref{app:closures} gives the equilibrium conditions and the code):
\begin{description}
\item[\texttt{sate}] The export activity of the modelled region sells into an export market with price $p^{PX}_{ir}$.
Equation \eqref{sate} is added as an auxiliary constraint, and the representative agent receives the foreign exchange value of the exports.
This implementation is particularly simple: it requires only a demand function for each export good and no explicit accounts of the rest of the world.
While listing the accounts of the rest of the world comes naturally in a single-country model derived from global \GTAP data, it means extra effort in single-country models calibrated to national \IOT{}s or \SAM{}s.
\item[\texttt{sate2}] The export demand elasticity is implemented as in the \LOE model of \citet{Schreiber2024}: the export activity combines the region's exports with a fixed factor owned by the rest of the world in a Cobb-Douglas technology to produce foreign exchange.
Setting the cost share of the fixed factor to $1/\sigma^m_i$ yields an export demand elasticity of $\sigma^m_i$ \citep[Proposition 2]{Schreiber2024}.
\end{description}
Both versions are included in the code to verify that they yield the same results.
They do: in all simulations of Section \ref{sec:comparison}, results of \texttt{sate} and \texttt{sate2} agree up to the six digits reported.
Below, I therefore only report results for \texttt{sate}, labelled \SATE.

\subsection{Model versions: \texttt{gtap} and \texttt{gtapingams}}\label{sec:versions}

The standard \GTAP model views domestic output as a homogeneous good that is sold domestically and exported alike ($\eta^{DX}_i=\infty$).
\citet{LanzRutherford2016} follow \citet{DeMeloRobinson1989} and introduce a finite elasticity of transformation between exported and domestically marketed output, which is needed to avoid corner solutions in the \SOE closure.
To keep the trade responses comparable to the standard \GTAP model, they compensate the finite transformation elasticity by a higher elasticity of substitution between imported and domestic goods (Appendix \ref{app:thetadx}): with parameter $\theta$ (\texttt{thetadx}), the elasticities are $\eta^{DX}_i=\sigma^D_i/(1-\theta)$ and $\sigma^{DM}_i=\sigma^D_i/\theta$, where $\sigma^D_i$ is the \GTAP elasticity of substitution between domestic and imported goods (\texttt{esubd}).
In line with \citet{LanzRutherford2016}, I use the same values in the \GMR and in the small open economy closures, and I consider two model versions:
\begin{description}
\item[\texttt{gtap}] $\theta=1$: $\eta^{DX}_i=\infty$ and $\sigma^{DM}_i=\sigma^D_i$, as in the standard \GTAP model.
\item[\texttt{gtapingams}] $\theta=0.5$, the default of \citet{LanzRutherford2016}: $\eta^{DX}_i=\sigma^{DM}_i=2\sigma^D_i$.
\end{description}
Table \ref{tab:elasticities} reports the resulting elasticities for the four sectors of the aggregation used below.

\begin{table}[t!]
\begin{center}
\caption{Trade and substitution elasticities of the \texttt{g20\_macro} aggregation (\GTAP 12, 2023).}\label{tab:elasticities}
\small
\begin{tabular}{lrrrrrr}
\toprule
 & & & & \multicolumn{2}{c}{\texttt{gtap}} & \texttt{gtapingams} \\
Sector & $\sigma^{VA}_i$ & $\sigma^D_i$ & $\sigma^m_i$ & $\sigma^{DM}_i$ & $\eta^{DX}_i$ & $\sigma^{DM}_i=\eta^{DX}_i$ \\
\midrule
Agriculture (\texttt{agr})   & 0.27 & 2.27 & 4.62 & 2.27 & $\infty$ & 4.53 \\
Manufacturing (\texttt{mfr}) & 1.18 & 3.37 & 7.14 & 3.37 & $\infty$ & 6.73 \\
Services (\texttt{ser})      & 1.37 & 1.91 & 3.80 & 1.91 & $\infty$ & 3.82 \\
Energy (\texttt{enr})        & 0.77 & 3.31 & 12.16 & 3.31 & $\infty$ & 6.63 \\
\bottomrule
\end{tabular}
\end{center}
\footnotesize{Notes: $\sigma^{VA}_i$: substitution between primary
factors (\texttt{esubva}); $\sigma^D_i$: \GTAP elasticity between domestic and imported goods (\texttt{esubd}); $\sigma^m_i$: substitution between import sources (\texttt{esubm}), which is also the export demand elasticity of the \SATE closure; $\sigma^{DM}_i$: elasticity between domestic and imported goods used in the model (\texttt{esubdm}); $\eta^{DX}_i$: transformation between domestic and export supply (\texttt{etadx}).
Values are aggregated from the \GTAP 12 values by the aggregation routine of \GTAPINGAMS.}
\end{table}

\subsection{Data: the \GTAP 12 Data Base and aggregations}\label{sec:data}

\GTAPtwelve reads the \GTAP 12 Data Base \citep{Aguiar2025}, which covers 163 regions, 65 sectors, and 8 primary factors and is available for seven reference years between 2004 and 2023.
All results below are based on the 2023 reference year.
As in \citet{LanzRutherford2016}, the data are filtered and rebalanced and then aggregated.
The regional aggregation \texttt{g20} of \citet{LanzRutherford2016} distinguishes the G20 economies and four composite regions (Table \ref{tab:regions}); it keeps all 65 sectors and 8 factors.
The aggregation \texttt{g20\_macro} has the same regions, four sectors (Table \ref{tab:sectors}) and four factors (labor, capital, land, and natural resources).
The simulations of Section \ref{sec:comparison} use the \texttt{g20\_macro} aggregation; the description of benchmark trade taxes in Appendix \ref{app:tariffs} uses the sectoral detail of the \texttt{g20} aggregation.

\begin{table}[t!]
\begin{center}
\caption{Regions of the \texttt{g20} and \texttt{g20\_macro} aggregations.}\label{tab:regions}
\small
\begin{tabular}{llll}
\toprule
Code & Region & Code & Region \\
\midrule
\texttt{arg} & Argentina & \texttt{jpn} & Japan \\
\texttt{anz} & Australia and New Zealand & \texttt{kor} & Korea \\
\texttt{bra} & Brazil & \texttt{mex} & Mexico \\
\texttt{can} & Canada & \texttt{rus} & Russia \\
\texttt{chn} & China and Hong Kong & \texttt{sau} & Saudi Arabia \\
\texttt{fra} & France & \texttt{zaf} & South Africa \\
\texttt{deu} & Germany & \texttt{tur} & Turkey \\
\texttt{gbr} & United Kingdom & \texttt{usa} & United States \\
\texttt{ind} & India & \texttt{reu} & Rest of the European Union \\
\texttt{idn} & Indonesia & \texttt{oex} & Other oil exporters \\
\texttt{ita} & Italy & \texttt{lic} & Other low income countries \\
 & & \texttt{mic} & Other middle income countries \\
\bottomrule
\end{tabular}
\end{center}
\end{table}

\begin{table}[t!]
\begin{center}
\caption{Sectors of the \texttt{g20\_macro} aggregation and the \GTAP 12 sectors they contain.}\label{tab:sectors}
\small
\begin{tabular}{lp{10.5cm}}
\toprule
Sector & \GTAP 12 sectors \\
\midrule
\texttt{agr} Agriculture & paddy rice, wheat, other cereal grains,
vegetables, fruit and nuts, oil seeds, sugar cane and beet, plant-based
fibres, other crops, cattle, other animal products, raw milk, wool,
forestry, fishing (14 sectors)\\
\texttt{enr} Energy & coal, crude oil, gas, petroleum and coal
products, electricity, gas manufacture and distribution (6 sectors)\\
\texttt{mfr} Manufacturing & other minerals, meat products, vegetable
oils, dairy, processed rice, sugar, other food, beverages and tobacco,
textiles, wearing apparel, leather, wood, paper, chemicals,
pharmaceuticals, rubber and plastics, other mineral products, ferrous
metals, other metals, metal products, motor vehicles, other transport
equipment, electronic equipment, computer, electronic and optical
products, machinery, other manufactures (27 sectors)\\
\texttt{ser} Services & construction, public administration and
defence, education, health, water, trade, accommodation and food
services, other transport, water and air transport, warehousing,
communication, financial services, insurance, real estate, business
services, recreational and other services, dwellings (18 sectors)\\
\bottomrule
\end{tabular}
\end{center}
\end{table}

\subsection{From \GTAPnine to \GTAPtwelve}\label{sec:update}

Table \ref{tab:changes} summarizes the changes of the code relative to the \GTAPnine distribution of \citet{LanzRutherford2016}.
The update to the \GTAP 12 Data Base mainly concerns the data routines: the \GDX version of the database replaces the \GEMPACK \HAR files, the reading routine has been rewritten accordingly, and the build stream is managed by the Python tool \texttt{doit}.
The model equations of \citet{LanzRutherford2016} are unchanged.
The extensions made for this paper add the \SATE closures, options to set $\theta$ and to scale $\sigma^m_i$ from the command line, and the update of average export taxes on exports to the rest of the world in the small open economy closures, which the original code held at benchmark values.

\begin{table}[t!]
\begin{center}
\caption{Changes relative to \GTAPnine \citep{LanzRutherford2016}.}\label{tab:changes}
\small
\begin{tabular}{lp{10.2cm}}
\toprule
File & Change \\
\midrule
\multicolumn{2}{l}{\textit{Update to the \GTAP 12 Data Base}} \\
\texttt{code/readgdx.gms} & Reads the \GDX distribution of \GTAP 12
(all reference years); replaces \texttt{flex2gdx.gms}. \\
\texttt{code/gtap12data.gms} & Successor of \texttt{gtap9data.gms};
model name \texttt{gtap12}, default reference year 2023. \\
\texttt{defines/*.map} & Aggregation maps for the 163 regions, 65
sectors and 8 factors of \GTAP 12. \\
\texttt{dodo.py} & Build stream (\texttt{doit}) for data processing,
checks, simulations, and figures (Table \ref{tab:dodo}). \\
\addlinespace
\multicolumn{2}{l}{\textit{Extensions for this paper}} \\
\texttt{code/mge.gms} & \SATE closures: commodities \texttt{PPX},
\texttt{PFIX}, auxiliary variables \texttt{EXPORTS}, \texttt{REVEXP},
new endowments of \texttt{RA(r)} and \texttt{ROW}, two variants of the
export activity \texttt{X(i,r)}; \texttt{thetadx} and \texttt{fsubm}
can be set from the command line. \\
\texttt{code/gtap12data.gms} & Flags \texttt{sate}, \texttt{sate2};
parameters \texttt{theta\_xi}, \texttt{fix0}; scaling of
$\sigma^m_i$ by \texttt{fsubm}. \\
\texttt{forensics/tgrid.gms} & Options \texttt{--mdl=sate|sate2},
\texttt{--thetadx}, \texttt{--fsubm}, \texttt{--soextax}; macro
\texttt{updrow} updating export taxes on exports to \textsc{row};
results for sectoral activity levels. \\
\texttt{forensics/*} & Batch runs (\texttt{ssasolve.gms}), merging of
results (\texttt{ssamerge.gms}), data for figures
(\texttt{tradetax.gms}, \texttt{process\_results.gms}). \\
\bottomrule
\end{tabular}
\end{center}
\end{table}

%	===================================================================
\section{Comparison of \SATE with the Global \GTAPINGAMS Model}\label{sec:comparison}
%	===================================================================

\subsection{Scenario design}\label{sec:design}

\paragraph{\bf Model versions.}
Each scenario is computed for the \texttt{gtap} and the \texttt{gtapingams} model versions of Section \ref{sec:versions}.

\paragraph{\bf Trade closures.}
I consider four trade closures:
\begin{description}
\item[\GMR] The standard global multiregional model, which serves as the reference.
\item[\GMRLS] The global model with the elasticities of substitution between imports from different sources ($\sigma^m_i$) halved.
\citet{Ahmad2021} show that different studies produce different estimates for Armington elasticities and that the \GTAP estimates are at the upper end of the range of estimates in the literature.
\GMRLS serves as a sensitivity analysis of the global model with respect to this uncertainty.
\item[\SOE] The small open economy model with fixed world prices of \citet{LanzRutherford2016}.
Without a finite elasticity of transformation between domestic and export supply, this closure is bound to produce bang-bang solutions in the \texttt{gtap} model version.
\item[\SATE] The small Armington trade economy (implementation \texttt{sate}), in which the world's demand for the exports of the modelled region has price elasticity $\sigma^m_i$ (benchmark values, not halved).
\end{description}

\paragraph{\bf Tariff scenarios.}
I reuse the tariff scenarios of \citet{LanzRutherford2016}.
For the region under consideration, all positive import tariff rates and all export tax rates on its exports (negative export subsidy rates) are scaled proportionally, such that the average trade tax rate (the value of these taxes relative to the value of the taxed trade flows at world prices) takes values from 0 to 10 percent in steps of one percentage point.
Trade taxes of all other regions remain at benchmark levels.
In the small open economy closures, the average export tax on exports to the rest of the world is updated accordingly.
Welfare is measured by the Hicksian equivalent variation of the region's representative agent, in percent of benchmark expenditure.
The benchmark average trade tax rate of each region is marked by a dot in the figures below; it is 2.24\% for the United States, 2.33\% for China, 4.83\% for South Africa, and 7.21\% for Argentina in the \texttt{g20\_macro} aggregation. %\footnote{%
  % The average is computed over the trade flows that carry positive taxes.
  % It therefore depends on the aggregation: in the more detailed \texttt{g20} aggregation, more tax-free flows drop out of the denominator, and the averages are 3.00\% for the United States and 3.52\% for China.
%}

Benchmark import tariffs and export subsidies are updated to the \GTAP 12 Data Base and the 2023 reference year.
The updated tariff rates are illustrated in Appendix \ref{app:tariffs}.

\subsection{Results}\label{sec:results}

\paragraph{\bf Welfare.}
Figures \ref{fig:W1} and \ref{fig:W2} show the welfare effects of the tariff scenarios for the four countries and both model versions.
Table \ref{tab:optimum} reports the tariff levels which maximize welfare in each case.
%, and Table \ref{tab:rmsdW} the root mean squared deviation (RMSD) of each closure from the \GMR results across the eleven tariff levels.

\begin{figure}[p!]
\centering
\input{jgea_new/figures/res_USA_W_gtap}\hfill
\input{jgea_new/figures/res_USA_W_gtapingams}

\vspace{0.1in}
\input{jgea_new/figures/res_CHN_W_gtap}\hfill
\input{jgea_new/figures/res_CHN_W_gtapingams}

\pgfplotslegendfromname{welfare_legend}
\caption{Welfare effect of unilateral changes of trade taxes, United States and China.
The dot marks the benchmark.
\SOE results outside the range shown are cut off.}\label{fig:W1}
\end{figure}

\begin{figure}[p!]
\centering
\input{jgea_new/figures/res_ZAF_W_gtap}\hfill
\input{jgea_new/figures/res_ZAF_W_gtapingams}

\vspace{0.1in}
\input{jgea_new/figures/res_ARG_W_gtap}\hfill
\input{jgea_new/figures/res_ARG_W_gtapingams}

\pgfplotslegendfromname{welfare_legend}
\caption{Welfare effect of unilateral changes of trade taxes, South Africa and Argentina.
The dot marks the benchmark.
\SOE results outside the range shown are cut off.}\label{fig:W2}
\end{figure}

\begin{table}[t!]
\begin{center}
\caption{Welfare-maximizing average trade tax rate (\%).}\label{tab:optimum}
\small
\input{jgea_new/figures/tab_optimum}
\end{center}
\footnotesize{Notes: Maximum over the grid of tariff levels 0, 1,
\ldots, 10\%, refined by quadratic interpolation around the best grid point.
``$\geq$10'': welfare is still increasing at the upper end of the grid.
``Benchmark'': benchmark average trade tax rate in the \texttt{g20\_macro} aggregation.}
\end{table}

% \begin{table}[t!]
% \begin{center}
% \caption{Welfare: range of \GMR results and deviations of the other closures from \GMR (percentage points of Hicksian EV).}\label{tab:rmsdW}
% \small
% \input{jgea_new/figures/tab_rmsd_W}
% \end{center}
% \footnotesize{Notes: Range: maximum minus minimum of the \GMR results
% over the eleven tariff levels.
% RMSD: root mean squared deviation from the \GMR results over the eleven tariff levels.
% Bold: \SATE deviates less from \GMR than \GMRLS.}
% \end{table}

The \SOE closure recommends free trade in seven of the eight cases (and an average rate of 1.4\% for Argentina in the \texttt{gtapingams} version).
This is the well-known consequence of fixed terms of trade: without market power, trade taxes only distort.
In the \texttt{gtap} version, where domestic and exported output are perfect substitutes, the \SOE welfare curve has a kink near the benchmark tariff level which reflects a bang-bang reallocation of output between sectors: for China, the welfare change ranges from $+4.0\%$ under free trade to $-4.0\%$ at an average rate of 10\%, compared with a range of less than one percentage point in the \GMR model.
The finite transformation elasticity of the \texttt{gtapingams} version tames these responses, but the \SOE welfare curves remain much steeper than the \GMR curves.

The \SATE closure follows the \GMR welfare curve closely.
%Its RMSD from the \GMR results is at most 0.19 percentage points, and it is below 0.1 percentage points in five of the eight cases, whereas the RMSD of the \SOE closure ranges from 0.16 to 3.12 percentage points.
For the United States and Argentina in the \texttt{gtap} version, the \SATE and \GMR curves are almost indistinguishable.
The welfare-maximizing tariff level, however, is not identified as well.
Welfare curves are flat around their maximum, the welfare gains relative to the benchmark are small (less than 0.1\% of expenditure in all \GMR and \SATE cases), and small differences in curvature shift the optimum considerably.
In the \texttt{gtap} version, the \SATE optimum coincides with the \GMR optimum for the United States (free trade) and differs from it by about 0.3 percentage points for China and Argentina.
In the \texttt{gtapingams} version, the differences are 0.7 percentage points for Argentina and 1.2 percentage points for the United States and China.
The largest deviations occur for South Africa, where the \SATE closure puts the optimum at 3.0--3.7\% instead of 5.6--5.7\%.
Still, the \SATE closure gets the order of magnitude of the optimal tariff right in all cases, which the \SOE closure does not.
\SATE also reproduces a pattern of the \GMR results that the \SOE closure cannot produce: Argentina and South Africa, whose benchmark protection is relatively high, would lose from unilateral free trade, whereas China and the United States would gain from it in the \texttt{gtap} version.

The remaining deviations between \SATE and \GMR should be put into perspective.
Halving the elasticities of substitution between import sources in the global model (\GMRLS) changes the results much more than switching from the \GMR to the \SATE closure.
%: the RMSD of \GMRLS from \GMR ranges from 0.14 to 0.90 percentage points, and the welfare-maximizing tariff level lies beyond the upper end of the grid (10\%) in all cases.
The intuition behind why substitution elasticities are so relevant to  locating the welfare optimum seems clear:
Lower Armington elasticities mean more market power of a country's exports in its export markets (and of other countries' exports in the country's import markets), so that higher import tariffs and export taxes improve the terms of trade without distorting quantities as much as with high elasticities.
Given that the the range of plausible values for  Armington elasticities is considerable \citep{Ahmad2021}, the uncertainty about the correct elasticities is a larger source of uncertainty about the welfare effects of trade policy than the approximation error of the \SATE closure.

\paragraph{\bf Sectoral activity levels.}
Figures \ref{fig:sec_large} and \ref{fig:sec_small} compare the changes of sectoral activity levels at an average trade tax rate of 10\%, where the differences between the closures are largest.
For each sector, the \GMR result is marked by a black dot on a vertical bar; the comparison to results of the \SATE, \GMRLS, and \SOE closures are illustrated by horizontal lines.
The scale of the x-axis is set by the \GMR, \SATE, and \GMRLS results; \SOE results outside this range are indicated by an arrow and their value.

\begin{figure}[p!]
\centering
\input{jgea_new/figures/dumbbell_g20_macro_USA_gtap}\hfill
\input{jgea_new/figures/dumbbell_g20_macro_USA_gtapingams}

\vspace{0.1in}
\input{jgea_new/figures/dumbbell_g20_macro_CHN_gtap}\hfill
\input{jgea_new/figures/dumbbell_g20_macro_CHN_gtapingams}

\pgfplotslegendfromname{dumbbell_legend}
\caption{Sectoral activity levels at an average trade tax rate of 10\% (\% change relative to the benchmark), United States and China.
\SOE results outside the range shown are indicated by arrows and their values.}\label{fig:sec_large}
\end{figure}

\begin{figure}[p!]
\centering
\input{jgea_new/figures/dumbbell_g20_macro_ZAF_gtap}\hfill
\input{jgea_new/figures/dumbbell_g20_macro_ZAF_gtapingams}

\vspace{0.1in}
\input{jgea_new/figures/dumbbell_g20_macro_ARG_gtap}\hfill
\input{jgea_new/figures/dumbbell_g20_macro_ARG_gtapingams}

\pgfplotslegendfromname{dumbbell_legend}
\caption{Sectoral activity levels at an average trade tax rate of 10\% (\% change relative to the benchmark), South Africa and Argentina.
\SOE results outside the range shown are indicated by arrows and their values.}\label{fig:sec_small}
\end{figure}

Across the 32 combinations of country, model version, and sector, the \SATE result is closer to the \GMR result than the \SOE result in 25 cases and closer than the \GMRLS result in 19 cases.
The \SOE closure fails most clearly in the \texttt{gtap} version, where domestic and exported output are perfect substitutes: an average trade tax rate of 10\% shrinks manufacturing by 15--36\% and expands services by 3--11\% in all four countries, whereas the \GMR model predicts changes of less than 1.5\% in these sectors.
In the \texttt{gtapingams} version, the \SOE results are of the right order of magnitude, but they still deviate substantially for energy, where they overstate the output response in the United States, China, and South Africa.

For the two small economies, the \SATE closure reproduces the \GMR results almost exactly (Figure \ref{fig:sec_small}).
For Argentina, its deviations do not exceed 0.03 percentage points in any sector and model version, and for South Africa they are at most 0.13 percentage points; in both countries they are mostly smaller than the deviations of \GMRLS.
The \GMRLS results are quite different in the energy sector of South Africa, where halving the Armington elasticities reverses the sign of the output change in both model versions.

The fit is weaker for the large economies (Figure \ref{fig:sec_large}).
The \SATE closure tends to overstate the output responses of the \GMR model: energy output in the United States declines by 9.7\% instead of 7.3\% (\texttt{gtap} version) and by 6.9\% instead of 5.8\% (\texttt{gtapingams} version); Chinese manufacturing declines by 1.6\% instead of 1.3\% (\texttt{gtap} version), and Chinese agriculture expands by 3.3\% instead of 2.9\% (\texttt{gtapingams} version).
In one case, US manufacturing in the \texttt{gtap} version, the \SATE closure gets the sign wrong: it predicts a decline of 0.5\% where the \GMR model predicts a small increase of 0.1\%.
For these economies, \GMRLS is often closer to the \GMR results than \SATE.
This is what one would expect: the United States and China are large in many of their export markets, so that the small-economy limit underlying the \SATE closure overstates the elasticity of export demand (Section \ref{sec:sate}), and the terms of trade of their imports, which are fixed in the \SATE closure, matter as well.

%	===================================================================
\section{Conclusion}\label{sec:concl}
%	===================================================================

Single-country \CGE models based on national \IOT{}s or \SAM{}s cannot represent terms of trade effects if they assume fixed world prices for both imports and exports.
This paper proposes the small Armington trade economy (\SATE) as a simple remedy: the world's demand for the country's exports has a constant price elasticity equal to the Armington elasticity of substitution between import sources, while import prices remain fixed.
The specification follows from the limit of an Armington world economy as the country becomes small \citep{CaliendoFeenstra2024}.
It needs no data about the rest of the world beyond the country's own exports and a set of Armington elasticities, and it is easy to implement: the version \texttt{sate} adds one price and two auxiliary constraints per export good.

The comparison with the global \GTAPINGAMS model, updated here to the \GTAP 12 Data Base, shows that the \SATE closure reproduces the welfare effects of unilateral changes of trade taxes far better than the \SOE closure, and that it gets the order of magnitude of the welfare-maximizing tariff level right, where the \SOE closure recommends free trade.
Its deviations from the global model are not larger than the uncertainty in the global model's results that is due to the uncertainty about the right Armington elasticities.
The fit of sectoral activity levels is mixed, and it is weakest for the largest economies, for which the small-economy limit is least appropriate.
The more elaborate calibration of a large open economy model by \citet{Schreiber2024}, which estimates export demand and import supply elasticities from the global model, can be expected to do better for such economies, at the price of requiring estimates of elasticities derived from an equivalent, global model.

While the \texttt{gtap} version of the model is closer to the original \GTAP model, the \texttt{gtapingams} version of \citet{LanzRutherford2016} is not without merit.
\citet{Soderbery2015} and \citet{Broda2008} estimate export supply elasticities, and the latter make them central to their argument for import tariffs based on the market power of importers.
The \GTAP model does not explicitly calibrate price elasticities of export supply, but export supply functions to a specific country are implied by market interactions, first within the exporting country and then among the different countries importing the commodity in question.
For a more explicit calibration of export supply functions, imperfect transformability between domestically provided and exported varieties can be of use.
For the experiments of this paper, the \SATE closure performs similarly well in both versions, whereas the \SOE closure is only usable in the \texttt{gtapingams} version.

\begin{acknowledgements}

This paper builds on the \GTAPINGAMS model and its documentation by Bruno Lanz and Thomas F.\ Rutherford \citep{LanzRutherford2016}.
Appendix \ref{app:model} reproduces large parts of their model documentation.
Any remaining errors are mine.

\end{acknowledgements}

\begin{appendices}
\section{Model Documentation}\label{app:model}

This appendix documents the \GTAPtwelve model for completeness.
Subsections \ref{sets_sec}--\ref{equil} and \ref{app:thetadx} reproduce, with minor edits, Sections 2, 3, and 4.3 of \citet{LanzRutherford2016}.
The edits concern (i) the dimensions of the \GTAP 12 Data Base, (ii) the choice of the elasticity of transformation between domestic and export supply in the closures used in this paper, and (iii) the valuation of the current account, which follows the \GAMS code.
Subsections \ref{app:closures} and \ref{app:programs} are new: they document the closures with respect to the rest of the world, including the \SATE closure proposed in this paper, and the programs.
The \MPSGE source code of the model is listed in Subsection \ref{app:code}.
Readers interested in the \LES and \CDE demand systems, which are part of the code but not used in this paper, are referred to Appendices C and D of \citet{LanzRutherford2016}.

\subsection{Notation}\label{sets_sec}

The notation used in the model is summarized in the Tables \ref{sets} - \ref{prices}.
Table \ref{sets} defines the various dimensions which characterize an instance of the model, including the set of sectors/commodities ($i,j$), the set of regions ($r,s)$, the set of factors of production ($f$).
Set $g$ combines the production sectors $i$ and private and public consumption demand (indices \texttt{"C"} and \texttt{"G"}) and investment demand (index \texttt{"I"}).
It allows for a much tighter formulation of the model as they can all be conceived of ``goods'' produced in similar fashion.
To simplify the exposition of the model, however, we describe private consumption, public consumption, and investment demand as stand alone components.

\renewcommand{\arraystretch}{1}
\renewcommand{\tabcolsep}{1mm}
\begin{table}[t!]
\begin{center}
\caption{Definitions of set indices.}
\label{sets}
\small
\begin{tabular}{lp{4in}}
\toprule
Set &Definition  \\
\addlinespace
\midrule
\addlinespace
$i,j$ & Sectors, an aggregation of the 65 sectors in the \GTAP 12 Data Base \\
$g$ & Production sectors $i$, plus private consumption \texttt{"C"},
public demand \texttt{"G"}, and investment \texttt{"I"} \\
$r,s$ & Regions, an aggregation of the 163 regions in the \GTAP 12 Data Base\\
$f$ & Factors of production (consisting of
\textit{mobile factors}, $f \in m_f$: five labor categories
(officials, managers, and professionals; technicians and associate
professionals; clerks; service and shop workers; agricultural and other
low-skilled workers) and capital, and \textit{sector-specific} (sluggish)
factors, $f \in s_f$: agricultural land and natural resources; in the
\texttt{g20\_macro} aggregation, the five labor categories are
aggregated to a single labor input)\\ \addlinespace \bottomrule
\end{tabular}
\end{center}
\vspace{0.4in}
\end{table}

Table \ref{activities} defines the primal variables (activity levels) which characterize an equilibrium.
The model determines values of all the variables except international capital flows, a parameter which would be determined endogenously in an intertemporal model.
Table \ref{activities} also displays the concordance between the variables and their \GAMS equivalents.

\renewcommand{\arraystretch}{1}
\renewcommand{\tabcolsep}{1mm}
\begin{table}[t!]
\begin{center}
\caption{Definitions of activity levels (quantity variables).}
\label{activities}
\small
\begin{tabular}{lllr}
\toprule
Variable &Definition  & GAMS variable & Benchmark (\GTAP) value \\
\addlinespace
\midrule
\addlinespace
$Y_{ir}$&       Production & \texttt{Y(i,r)} &\texttt{vom(i,r)}\\
\\
$C_{r}$&        Discretionary consumption & \texttt{Y("c",r)} &\texttt{vom("c",r)}\\
\\
$G_{r}$&        Aggregate public & \texttt{Y("g",r)}&\texttt{vom("g",r)} \\
\\
$I_{r}$&        Aggregate investment & \texttt{Y("i",r)}&\texttt{vom("i",r)} \\
\\
$M_{ir}$&       Aggregate imports & \texttt{M(i,r)}&\texttt{vim(i,r)}\\
\\
$X_{ir}$&       Trade flows with rest of world regions &
\texttt{X(i,r)}&\texttt{vem(i,r)}\\
\\
$\textit{FT}_{fr}$&      Factor transformation &
\texttt{FT(f,r)}&\texttt{evom(f,r)}\\
\\
$\textit{YT}_j$  &       International transport services & \texttt{YT(j)}&\texttt{vtw(j)}\\
\addlinespace
\bottomrule
\end{tabular}
\end{center}
\vspace{0.4in}
\end{table}

Table \ref{prices} defines the relative price variables for goods and factors in the model.
As is the case in any Shoven-Whalley \CGE model, the equilibrium conditions determine \textit{relative} rather than \textit{nominal} prices.

\renewcommand{\arraystretch}{1}
\renewcommand{\tabcolsep}{1mm}
\begin{table}[t!]
\begin{center}
\caption{Definitions of price variables.}
\label{prices}
\small
\begin{tabular}{lp{3in}r}
\toprule
Variable &Definition & GAMS variable\\
\addlinespace
\midrule
\addlinespace
$p^C_{r}$      &       Final demand price index for the Cobb-Douglas
demand system &\texttt{P("C",r)}\\
\\
$p^G_{r}$      &       Public provision price index  &\texttt{P("G",r)}\\
\\
$p^I_{r}$      &       Investment price index &\texttt{P("I",r)}\\
\\
$p^Y_{ir}$     &       Supply price, gross of indirect producer
taxes (when $\eta^{DX}_{i} < \infty$, the supply price is
replaced by and export price, $p^X_{ir}$ and domestic supply price,
$p^D_{ir}$) &\texttt{P(i,r)}\\
\\
$p^M_{ir}$     &       Import price, gross of export
taxes and tariffs.&\texttt{PM(i,r)}\\
\\
$p^E_{ir}$     &       Export price from endogenous or exogenous
regions.&\texttt{PE(i,r)}\\
\\
$p^T_j$         &       Marginal cost of transport services
&\texttt{PT(j)}\\
\\
$p^F_{f,r}$     &       Price of mobile primary factors ($m_f$
include labor, land, and resources) &\texttt{PF(f,r)}\\
\\
$p^S_{s_f,ir}$     &       Price of sector-specific primary
factors&\texttt{PS(sf,i,j)}\\
\addlinespace
\bottomrule
\end{tabular}
\end{center}
\vspace{0.4in}
\end{table}

While the core \GTAPINGAMS model is base on Cobb-Douglas final demand, the canonical models incorporate logic for two additional demand systems: the \LES and the \CDE demand system.
In particular, the \LES representation distinguishes between subsistence demand and discretionary demand, and we define the additional variables in Table \ref{dsysprices}.
Similarly, we define an Armington composite price index for products entering \CDE final demand.
(Details of these demand systems are provided in Appendices C and D of \citet{LanzRutherford2016}.)

\begin{table}[t!]
\begin{center}
\caption{Additional variables for \LES and \CDE demand systems.}
\label{dsysprices}
\small
\begin{tabular}{lp{3in}r}
\toprule
Variable &Definition & GAMS variable\\
\addlinespace
\midrule
\addlinespace
$SD_{r}$&       Subsistence demand & \texttt{Y("sd",r)} \\
\\
$DD_{r}$&       Discretionary demand & \texttt{Y("dd",r)} \\
\\
$p^{DD}_{r}$   &       Discretionary demand price index &\texttt{P("dd",r)}\\
\\
$p^{SD}_{r}$   &       Subsistence demand price index &\texttt{P("sd",r)}\\
\\
$p^A_{iCr}$     &      Armington composite price representing the
market price of goods entering \CDE demand.&\texttt{PA(i,r)}\\
\addlinespace
\bottomrule
\end{tabular}
\end{center}
\vspace{0.4in}
\end{table}

Finally, Table \ref{taxes} reports the definition of tax and subsidy rates applied in the model, both in terms of the notation employed to describe the model and that used in the \GAMS code.
Note that revenues from taxes and subsidy expenditures do not appear as explicit variables in the \GTAP Data Base and are defined on the basis of expenditures and tax rates.
We come back to this below.

\renewcommand{\arraystretch}{1}
\renewcommand{\tabcolsep}{1mm}
\begin{table}[t!]
\begin{center}
\caption{Tax and subsidy rates (net basis unless noted).}
\label{taxes}
\small
\begin{tabular}{lp{6cm}lr}
\toprule
Parameter & Definition &  & \GAMS Parameter \\
\addlinespace
\midrule
\addlinespace
$t^o_{ir}$ & Output taxes (gross basis)& & \texttt{rto(i,r)}    \\
\\
$t^{f}_{fir}$   &Factor taxes &  & \texttt{rtf(f,i,r)}   \\
\\
$t^{fd}_{ijr}$ & Intermediate input taxes &Domestic&\texttt{rtfd(i,j,r)}      \\
$t^{fi}_{ijr}$ &                          &Imported& \texttt{rtfi(i,j,r)}      \\
\\
$t^{pd}_{ir}$ & Consumption taxes  & Domestic& \texttt{rtfd(i,"C",r)}     \\
$t^{pi}_{ir}$ &                 &Imported& \texttt{rtfi(i,"C",r)} \\
\\
$t^{gd}_{ir}$ & Public demand taxes &Domestic & \texttt{rtfd(i,"G",r)}   \\
$t^{gi}_{ir}$&                 &Imported & \texttt{rtfi(i,"G",r)} \\
\\
$t^{gd}_{ir}$ & Investment demand taxes &Domestic & \texttt{rtfd(i,"I",r)}   \\
$t^{gi}_{ir}$&                 &Imported& \texttt{rtfi(i,"I",r)} \\
\\
$t^{xs}_{isr}$  &Export subsidies&  & \texttt{rtxs(i,s,r)}   \\
\\
$t^{ms}_{isr}$  &Import tariffs& & \texttt{rtms(i,s,r)}\\
\addlinespace
\bottomrule
\end{tabular}
\end{center}
\vspace{0.4in}
\end{table}

\subsection[Accounting identities]{Benchmark data structure and accounting identities}

The economic structure underlying the \GTAP dataset and model is illustrated in Figure \ref{fig:region}.
Symbols in this flow chart correspond to variables in the economic model (see Table \ref{activities}): $Y_{ir}$ is the production of good $i$ in region $r$, $C_r$, $I_r$, and $G_r$ portray private consumption, investment, and public demand, respectively, $M_{ir}$ portrays the import of good $i$ into region $r$, $\textit{RA}_r$ stand for representative consumers, and $FT_{s_{f}r}$ is the activity through which the set of sector-specific factors of production ($s_f$) are allocated to individual sectors.
Further, solid lines represent commodity and factor market flows, while dotted lines indicate tax revenues and transfers.

\begin{sidewaysfigure}
\begin{center}
\includegraphics{jgea_new/figures/gtapflowchart.tikz}
\caption{Regional economic structure.}
\label{fig:region}
\end{center}
\end{sidewaysfigure}

Domestic and imported goods markets are represented by horizontal lines at the top of the figure.
Domestic production ($\textit{vom}_{ir}$) is distributed to exports ($\textit{vxmd}_{irs}$), international transportation services ($\textit{vst}_{ir}$), intermediate demand ($\textit{vdfm}_{ijr}$), household consumption ($\textit{vdfm}_{iCr}$), investment ($\textit{vdfm}_{iIr}$), and government consumption ($\textit{vdfm}_{iGr}$).\footnote{Recall that in the \GAMS implementation of the model the index \texttt{"g"} includes all sectors represented in the model plus private consumption \texttt{"C"}, public demand \texttt{"G"}, and investment \texttt{"I"}.
(See Table \ref{sets}.)
For the LES demand representation, it also includes price indices for discretionary demand \texttt{"dd"} and subsistence demand \texttt{"sd"}. } The accounting identity in the \GTAP dataset is thus:
$$\textit{vom}_{ir} = \sum_s \textit{vxmd}_{irs} + \textit{vst}_{ir}
+ \sum_j \textit{vdfm}_{ijr} + \textit{vdfm}_{iCr} + \textit{vdfm}_{iIr} + \textit{vdfm}_{iGr}\,,$$ where $j$ indexes all goods.
Similarly, imported goods (with aggregate value $\textit{vim}_{ir}$) enter intermediate demand ($\textit{vifm}_{ijr}$), private consumption ($\textit{vifm}_{iCr} + \textit{vifm}_{i\textit{SD}r}$), and public consumption ($\textit{vifm}_{iGr}$).
The accounting identity for these flows is thus:
$$\textit{vim}_{ir} = \sum_j \textit{vifm}_{ijr} +
\textit{vifm}_{iCr} + \textit{vifm}_{i\textit{SD}r} + \textit{vifm}_{iGr}\,.$$

Inputs to production of good $i$ ($Y_{ir}$) include intermediate inputs (domestic $\textit{vdfm}_{ijr}$ and imported $\textit{vifm}_{ijr}$), mobile factors of production ($\textit{vfm}_{m_{f},ir}$, where $m_{f}$ is a subset of the set $f$ designating all factors of production), and sector-specific factors of production ($\textit{vfm}_{s_{f},ir}, s_f \subset f$).
Factor market equilibrium is given by an identity relating the value of factor payments to factor income:
$$\sum_i \textit{vfm}_{fir} = \textit{evom}_{fr}\,,$$ and factor
earnings accrue to households.

International market clearance conditions require that region $r$ exports of good $i$ ($\textit{vxmd}_{ir}$ at the top of the figure) equal the imports of the same good from the same region summed across all trading partners ($\textit{vxmd}_{isr}$ at the bottom of the figure): $$\textit{vxm}_{ir} = \sum_s \textit{vxmd}_{irs}\,,$$ where $s$, an alias for $r$, indexes regions.
Likewise, market clearance conditions apply for international transportation services.
The supply-demand balance in the market for transportation service $j$ requires that the sum across all regions of service exports ($\textit{vst}_{ir}$, at the top of the figure) equals the sum across all bilateral trade flows of service inputs ($\textit{vtwr}_{jisr}$ at the bottom of the figure):
$$\sum_r \textit{vst}_{jr} = \sum_{isr} \textit{vtwr}_{jisr}$$

Turning to tax revenues and transfers, shown as dotted lines in figure \ref{fig:region}, flows labeled with $\mathcal{R}$ correspond to tax revenues.
For each country, tax flows consist of indirect taxes on production/exports of each good ($\mathcal{R}^Y_{ir}$), on consumption ($\mathcal{R}^C_r$), on public demand ($\mathcal{R}^G_r$), and on imports ($\mathcal{R}^M_{ir}$).
The regional budget constraint thus relates tax payments ($\mathcal{R}^Y_{ir}, \mathcal{R}^C_r, \mathcal{R}^G_r, \mathcal{R}^M_{ir}$), factor income ($\textit{evom}_{fr}$), and the current account deficit (i.e., net transfers from abroad, $\textit{vb}_r$) to total private consumption expenditure $\textit{vom}_{Cr}$, total public consumption expenditure $\textit{vom}_{Gr}$, and total investment $\textit{vom}_{Ir}$, yielding:
$$\textit{vom}_{Cr} + \textit{vom}_{Gr} + \textit{vom}_{Ir} = \sum_f
\textit{evom}_{fr} + \sum_i \mathcal{R}^Y_{ir} + \mathcal{R}^C_r + \mathcal{R}^G_r + \sum_i \mathcal{R}^M_{ir} + \textit{vb}_r\,.
$$

To this point we have outlined two types of consistency conditions which are part of any \GTAP Data Base: market clearance (supply = demand for all goods and factors), and income balance (net income = net expenditure).
A third set of identities involve net operating profits by all sectors in the economy.
In the core \GTAP model ``production'' takes place under conditions of perfect competition with constant returns to scale, hence there are no excess profits, and the cost of inputs must equal the value of outputs.
This condition applies for each production sector:

\begin{itemize}
\item[$Y_{ir}$:] $\sum_f \textit{vfm}_{fir} + \sum_j \left(\textit{vdfm}_{ijr} + \textit{vifm}_{ijr} \right) + \mathcal{R}^Y_{ir} = \textit{vom}_{ir}\,,$
\item[$C_{r}$:] $\sum_i \left(\textit{vdfm}_{iCr} + \textit{vifm}_{iCr} \right) + \mathcal{R}^C_{ir} = \textit{vom}_{Cr}\,,$
\item[$I_{r}$:] $\sum_i \textit{vdfm}_{iIr} = \textit{vom}_{Ir}$
\item[$G_{r}$:] $\sum_i \left(\textit{vdfm}_{iGr} + \textit{vifm}_{iGr} \right) + \mathcal{R}^G_{ir} = \textit{vom}_{Gr}\,,$
\item[$M_{ir}$:] $\sum_s \left(\textit{vxmd}_{isr} + \sum_j \textit{vtwr}_{jisr} \right) + \mathcal{R}^M_{ir} = \textit{vim}_{ir}$
\item[$FT_{fr}$:] $\textit{evom}_{fr} = \sum_i \textit{vfm}_{fir}$
\end{itemize}

\subsection[Optimization problems]{Decentralized optimization problems}

The benchmark identities presented in the previous section indicate the market clearance, zero profit, and income balance conditions which define the \GTAP model.
The displayed equations do not, however, characterize the behavior of agents in the model.
In a competitive equilibrium setting, the standard assumption of optimizing atomistic agents applies for both producers and consumers.
This section lays out the optimization problem of each component in the model, and thereby provides the structure of production technology (production functions) and preferences (characterizing final demand), as well as the representation of trade.
The treatment of the rest of the world, which distinguishes the \GMR model from the small open economy closures, is described in Appendix \ref{app:closures}.

Note that in order to simplify notation, we denote decision variables corresponding to the benchmark data structures with the initial ``v'' replaced by ``d.''
Hence, while $vdfm_{jir}$ represents benchmark data on intermediate demand for good $j$ in the production of good $i$ in region $r$, $ddfm_{jir}$ represents the corresponding decision variable in the equilibrium model.
This approach to the scaling of variables is consistent with the \GAMS code, and it provides a flexible and transparent approach with respect to the calibration of activity variables.

\subsubsection{Production technology}

Starting with producers, profit maximization in the constant returns to scale setting is equivalent to cost minimization subject to technical constraints.
For sector $Y_{ir}$ we characterize input choices as though they arose from minimization of unit costs:
\begin{eqnarray}\label{profitmax}
\min_{\textit{ddfm,dfm,difm}} \quad & c^D_{ir} + c^M_{ir} +
c^F_{ir} & \\
&s.t.&c^D_{ir} = \sum_j p^Y_{jr} (1 + t^{fd}_{jir})
\textit{ddfm}_{jir} \nonumber\\
&&c^M_{ir} = \sum_j p^M_{jr} (1 + t^{fi}_{jir}) \textit{difm}_{jir}\nonumber\\
&&c^F_{ir} = \sum_f (p^F_{fr}|_{f \in m_f} + p^S_{fir}|_{f \in s_f}) (1 +
t^f_{fir} ) \textit{dfm}_{fir}\nonumber\\
&&F_{ir}(\textit{ddfm,difm,dfm}) = Y_{ir}\nonumber
\end{eqnarray}
where $F(\cdot)$ represents the production function, which is described by a nested constant-elasticity-of-substitution (CES) form, with structure displayed in Figure \ref{prodfunY}.

\begin{figure}[t!]
\centering
\includegraphics{jgea_new/figures/prodfunY.tikz}
\vspace{0.1in}
\caption{CES nesting structure for production function $Y_{ir}=F_{ir}(\textit{ddfm,difm,dfm})$.}
\label{prodfunY}
\vspace{0.3in}
\end{figure}

In the figure, $\sigma$ values in different nests represent substitution elasticities between inputs, with $\sigma^D_i = \textit{esubd}_i$ measuring substitution possibility between intermediate inputs produced domestically and imported from abroad which are similarly a composite of imports from varieties from different regions.
$\textit{esubva}_i$ represents the elasticity of substitution between primary inputs in the value added nest.
(Both of these parameters are provided in the \GTAP Data Base \citep[see][]{HertelvanderMensbrugghe2015}.)
Note further that the specific source of tax revenue is indicated in this figure, consisting of output taxes, taxes on intermediate inputs, and taxes on factor demands, all of which are applied on an ad-valorem basis.

One important modelling choice concerns the differentiation of output for domestic and export markets.
In particular, the supply of goods to domestic and export markets are portrayed as arising from the following profit-maximization problem:
\begin{eqnarray}\label{cet_dx}
\max_{\textit{dxm},\textit{ddm}} \quad && p^D_{ir} \textit{ddm}_{ir}
+ p^X_{ir} \textit{dxm}_{ir}
\\
&s.t.& G_{ir}(\textit{ddm}_{ir},\textit{dxm}_{ir}) = Y_{ir}
\end{eqnarray}
where $G$ is the \CET function with structure illustrated in Figure \ref{prodfunDX}.
In the standard \GTAP model, transformation elasticities $\eta^{DX}_{ir} = \textit{etrndx}_{i}$ are infinite.
This implies that the supply price of output $p^Y_{jr}$ is the same if output is used locally or in a different region.
With $\textit{etrndx}_{i} < \infty$, the output price may differ when supplied domestically or abroad.
In this case, $p^Y_{jr}$ in equation \eqref{profitmax} is replaced by $p^D_{jr}$.
In \GTAPINGAMS, the choice of $\textit{etrndx}_{i}$ is governed by the parameter \texttt{thetadx} (see Appendix \ref{app:thetadx}) and applies to both the \GMR and the small open economy closures: \texttt{thetadx}$=1$ yields $\textit{etrndx}_{i}=\infty$ (the \texttt{gtap} model version of this paper), \texttt{thetadx}$=0.5$ yields a finite elasticity (the \texttt{gtapingams} model version, the default of \citealp{LanzRutherford2016}).

\begin{figure}[t!]
\centering
\includegraphics{jgea_new/figures/prodfunDX.tikz}
\vspace{0.1in}
\caption{CET transformation between domestic and export markets $\textit{Y}_{ir} = G_{ir}(\textit{dxm},\textit{ddm})$.}
\label{prodfunDX}
\vspace{0.3in}
\end{figure}

\subsubsection{Preferences and final demand}

Private consumption consistent with utility maximization is portrayed by minimization of the cost of a given level of aggregate consumption:
\begin{eqnarray}\label{utils}
\min_{\textit{ddfm}_{iCr},\textit{difm}_{iCr}} \quad &&\sum_i p^Y_{ir}
(1 + t^{pd}_{ir}) \textit{ddfm}_{iCr} + p^M_{ir} (1 +t^{pi}_{ir})
\textit{difm}_{iCr} \\
&s.t.& H_r(\textit{ddfm}_{iCr},\textit{difm}_{iCr}) = C_{ir} \nonumber
\end{eqnarray}
where $H_r$ represents final demand from the representative consumer.

Final demand in the core model is characterized by Cobb-Douglas preferences.
Alternative specifications included in the model logic include both the \LES and \CDE expenditure.\footnote{The \LES model involves subsistence and discretionary expenditure, which are represented through composite price indices representing \CES aggregates of both domestic and imported inputs, \texttt{P("DD",r)} and \texttt{P("SD",r)}, respectively.
The \CDE model requires the introduction of trade aggregation activities (\texttt{A(i,r)}) and associated Armington composite price indices, \texttt{PA(i,r)}.}
The nested discretionary and subsistence demand functions are is displayed in Figure \ref{prodfunC}.

\begin{figure}[t!]
\centering
\includegraphics{jgea_new/figures/prodfunC.tikz}
\vspace{0.1in}
\caption{LES nesting for discretionary and subsistence consumption.}
\label{prodfunC}
\vspace{0.3in}
\end{figure}

\subsubsection{Government and public consumption}

Public consumption in the model is represented as a fixed coefficient (Leontief) aggregation of domestic-import composites.
This formulation introduces substitution at the second level between domestic and imported inputs while holding sectoral commodity aggregates constant.
Figure \ref{prodfunG} illustrates the functional form.

\begin{figure}[t!]
\centering
\includegraphics{jgea_new/figures/prodfunG.tikz}
\vspace{0.1in}
\caption{Nesting structure for public consumption $G_r = G_r(\textit{ddfm}_{iGr},\textit{difm}_{iGr})$.}
\label{prodfunG}
\vspace{0.3in}
\end{figure}

\subsubsection{International trade}

The choice among imports from different trading partners is based on Armington's idea of regionally differentiated products.
The following cost minimization problem formalizes this choice:
\begin{eqnarray}\label{optA}
\min_{\textit{dxmd,dtwr}} \quad &&\sum_s (1+t^{ms}_{isr}) \left(p^Y_{is}
(1-t^{xs}_{isr}) \textit{dxmd}_{isr} +
\sum_j p^T_j \textit{dtwr}_{jisr} \right)\\
&s.t.&A_{ir}(\textit{dxmd,dtwr}) = M_{ir} \nonumber
\end{eqnarray}
where $A$ is the import aggregation function, described by the nested \CES-Leontief function shown in Figure \ref{prodfunA}.
When $\textit{etrndx}_{i}<\infty$, $p^Y_{is}$ in (\ref{optA}) is replaced by $p^X_{is}$.

\begin{figure}[t!]
\centering
\includegraphics{jgea_new/figures/prodfunA.tikz}
\vspace{0.1in}
\caption{Armington aggregation of traded goods $M_{ir} = A_{ir}(\textit{dxmd,dtwr})$.}
\label{prodfunA}
\vspace{0.3in}
\end{figure}

Note that transportation services enter on a proportional basis with imports from different countries, reflecting differences in unit transportation margins across different goods and trading partners.
Therefore, substitution at the top level in an Armington composite involves trading off of both imported goods and associated transportation services.
Trade flows are subject to export subsidies and import tariffs, with subsidies paid by government in the exporting region, and tariffs collected by government in the importing region.

The provision of international transportation services is modeled through an aggregation of transportation services exported from countries throughout the world.
More specifically, we consider the following cost minimization problem for the aggregation of transportation services:
$$\min_{\textit{dst}} \quad \sum_r p^Y_{ir} \textit{dst}_{ir} \quad
\text{s.t.}
\quad T_i(\textit{dst}) = \textit{YT}_i$$ where the aggregation function $T_i$ combines transport service exports from multiple regions.
The functional form which aggregates services from different regions is illustrated in Figure \ref{prodfunYT}.

\begin{figure}[t!]
\centering
\includegraphics{jgea_new/figures/prodfunYT.tikz}
\vspace{0.1in}
\caption{International transportation services aggregator $\textit{YT}_j = T_j(\textit{dst})$.}
\label{prodfunYT}
\vspace{0.3in}
\end{figure}

\subsubsection{Supply of sector-specific factors}

Land and natural resources are portrayed as sector-specific factors of production.
These are supplied through constant-elasticity-of-transformation (CET) production function which allocates composite factors to sectoral markets.
Formally, the supply of sectoral factors of production is modeled through the following profit-maximization problem:
\begin{eqnarray}\label{cet_input}
\max_{\textit{dfm}} \quad && \sum_j \textit{dfm}_{s_{f},jr} p^S_{s_{f},jr}\\
&s.t.& \Gamma_{s_{f},r}(\textit{dfm}) = \textit{evom}_{s_{f},r}
\end{eqnarray}
where $\Gamma$ is the \CET function with structure illustrated in Figure \ref{prodfunFT}.
Note that in the figure $\eta$ represents transformation elasticities provided in the \GTAP Data Base.

\begin{figure}[t!]
\centering
\includegraphics{jgea_new/figures/prodfunFT.tikz}
\vspace{0.1in}
\caption{Sector-specific factor \CET transformation function $\textit{FT}_{s_{f}} = \Gamma_{s_{f}}(\textit{dfm})$.}
\label{prodfunFT}
\vspace{0.3in}
\end{figure}

\subsection[Equilibrium conditions]{Equilibrium Conditions}\label{equil}

An Arrow-Debreu model concerns itself with the interactions of decentralized decisions by consumers and producers in markets.
\citet{Lars} proposed a representation of this class of models in which two types of equations define an equilibrium: zero profit and market clearance.
The corresponding variables defining an equilibrium are activity levels (for constant-returns-to-scale firms) and commodity prices.\footnote{Under a maintained assumption of perfect competition, \citet{Lars} may characterize technology as constant-returns-to-scale without loss of generality.
Specifically, decreasing returns are accommodated through introduction of a specific factor, while increasing returns are inconsistent with the assumption of perfect competition.
Note that in this environment zero excess profit is consistent with free entry for atomistic firms producing an identical product.}
Here we extend Mathiesen's framework with a third class of variables corresponding to consumer income levels.
Commodity markets encompass primary endowments of households with producer outputs.
In equilibrium the aggregate supply of each good must be at least as great as total intermediate and final demand.
Initial endowments are exogenous.
Producer supplies and demands are defined by producer activity levels and relative prices.
Final demands are determined by market prices.

Economists who have worked with conventional textbook equilibrium models can find Mathiesen's framework to be somewhat opaque because many quantity variables need not be explicitly specified in the model.
Variables such as final demand by consumers, factor demands by producers, and commodity supplies by producers, are defined implicitly in Mathiesen's model.
For example, given equilibrium prices for primary factors, consumer incomes can be computed, and given income and goods prices, consumers' demands can then be determined.
The consumer demand functions are written down in order to define an equilibrium, but quantities demanded need not appear in the model as separate variables.
The same is true of inputs or outputs from the production process: relative prices determine conditional demand, and conditional demand times the activity level represents market demand.
Omitting decisions variables and suppressing definitional equations corresponding to intermediate and final demand provides significant computational advantages at the cost of a somewhat more complicated model statement.

In the following, we detail (i) zero profit conditions, (ii) market clearance conditions, and (iii) income balance conditions, which in the present case is equivalent to the regional budget constraint.
These three sets of conditions form the basic system of equation to be solved.
Note that the actual code for the model is implemented both in the algebraic mixed complementarity format \citep[\GAMS/\MCP, see][]{Rutherford1995} and through the more compact formulation afforded by the Mathematical Programming System for General Equilibruim (\MPSGE) syntax \citep{Rutherford1999}.

\subsubsection[Zero profit]{Zero profit (arbitrage) conditions}

All production activities in the model are represented by constant-returns-to-scale technologies, and markets are assumed to operate competitively with free entry and exit.
As a consequence, equilibrium profits are driven to zero and the price of output reflects the cost of inputs.
The following sets of equations relating output price to marginal cost are part of the definition of an equilibrium.\footnote{To retain consistency with the \MCP format, we express zero profit conditions as ``oriented equations,'' with marginal (=average) cost on the \LHS and marginal (=average) revenue on the \RHS.}

The calculation of unit cost and unit revenue functions involves the definition of a number of ancillary variables (that do not appear in the \GAMS code as explicit choice variables).
In the following we define ancillary variables in un-numbered equations, indicating that these variables are ``optional'' in the sense that they may be substituted out of the non-linear system of equations.
Moreover, we use the symbol $\theta$ to portray value shares from the base year data.
In most cases subscripts on these value shares are omitted in order to economize on notation.
Finally, to denote benchmark values we use an overline, so that $\overline{t}^{pd}_{ir}$ represents the benchmark value of ${t}^{pd}_{ir}$.

\paragraph{\bf Sectoral production (\texttt{Y(j,r)})}

Sectoral production combines intermediate inputs with a value-added nest combining primary inputs (see Figure \ref{prodfunY}).
The unit cost of value-added is a \CES composite of skilled and unskilled labor, land, resources, and capital inputs to production, gross of taxes.
Factor inputs may be sector-specific or mobile across sectors:
$$p^{pf}_{fjr} = \left\{\begin{array}{lr}
p^F_{fr} \frac{(1+t^f_{fjr})}{1+\overline{t}^f_{fjr}} & f \in m_f\\
p^S_{fjr} \frac{(1+t^f_{fjr})}{1+\overline{t}^f_{fjr}} & f \in s_f
\end{array} \right.$$
and the unit cost function is given by:
$$c^f_{jr} = \left( \sum_f \theta_f \left( p^{pf}_{fjr}
\right)^{1-\sigma} \right)^{1/(1-\sigma)},$$ in which $\theta_f$ represents the factor share of value added [{\small \texttt{theta\_vfm(f,g,r)}}].

The user cost of intermediate inputs differs from the market price due to the presence of taxes on intermediate inputs:
$$p^d_{ijr} = p^Y_{ir} \frac{1+t^{fd}_{ijr}}{1+\overline{t}^{fd}_{ijr}}$$
$$p^i_{ijr} = p^M_{ir}
\frac{1+t^{fi}_{ijr}}{1+\overline{t}^{fi}_{ijr}}$$ A \CES cost function describes the minimum cost of a bundle of domestic and imported inputs to production, based on benchmark value shares and an elasticity of substitution $\sigma [\texttt{esubd(i)}]$:
$$c^i_{ijr} = \left( \theta (p^d_{ijr})^{1-\sigma} +
             (1-\theta) (p^i_{ijr})^{1-\sigma} \right)^{1/(1-\sigma)},$$ in which $\theta$ represents the domestic share of the Armington composite \hfill \break [{\small \texttt{theta\_vdfm(i,j,r)}}].

Unit cost of sectoral output is then a Leontief (linear) composite of the costs of intermediate and value-added composite inputs, based on base-year value shares:
$$c^Y_{jr} = \sum_i \theta_i c^i_{ijr} + \theta^f c^f_{jr},$$
in which $\theta_i$ represents the cost share of intermediate input $i$ [{\small \texttt{theta\_cm0(i,j,r)}}]; and $\theta^f$ represents the cost share of value added in sectoral output [{\small \texttt{theta\_cf0(j,r)}}].

Having formulated the unit cost function, it is possible to compactly portray the zero profit condition for $y_{jr}$.
In equilibrium, the marginal cost of supply equals the market price, net of taxes:
\begin{equation}
c^Y_{jr} = p^Y_{jr} \frac{1-t^o_{jr}}{1-\overline{t}^o_{jr}}
\end{equation}
When $\textit{etrndx}_{i}<\infty$, production for domestic and export markets is differentiated, and we replace $p^Y_{jr}$ by a unit revenue function,
$$r^Y_{jr} = \left( \theta (p^D_{jr})^{1+\eta} +
(1-\theta) (p^X_{jr})^{1+\eta} \right)^{1/(1+\eta)},$$ in which $\theta$ represents the export share of output [{\small \texttt{theta\_vxm(j,r)}}].

\paragraph{\bf Consumer demand (\texttt{Y("c",r)})}

In previous versions of \GTAPINGAMS the consumer price index represented a Cobb-Douglas demand system with an Armington substitution between domestic and imported goods prices gross of tax.
This model is retained as one option in the new model, price indices for domestic and imported goods are given by:
$$p^{dc}_{ir} = p^Y_{ir} \frac{1+t^{pd}_{ir}}{1+\overline{t}^{pd}_{ir}},$$
and
$$p^{ic}_{ir} = p^M_{ir} \frac{1+t^{pi}_{ir}}{1+\overline{t}^{pi}_{ir}},$$
and the Armington composite price of good $i$ is a \CES composite price defined over these indices:
$$p^c_{ir} = \left(\theta (p^{dc}_{ir})^{1-\sigma} +
                (1-\theta) (p^{ic}_{ir})^{1-\sigma} \right)^{1/(1-\sigma)},$$ in which $\theta$ represents the domestic demand share [{\small \texttt{theta\_vdfm(i,"c",r)}}].

In addition to the Cobb-Douglas representation of the final demand system, our core model includes a \LES, which combines a Cobb-Douglas discretionary demand with a Leontief subsistence demand activity.
The added complexity of the \LES specification permits the model to be calibrated both to the income and \textit{average} price elasticities of demand.
(See Appendix C of \citet{LanzRutherford2016} for details.)
The price indices for discretionary and subsistence demand are defined by:
\begin{equation}
p^{DD}_r = \prod_i (p^c_{ir})^{\theta_i},
\end{equation}
in which $\theta_i$ represents the discretionary value share [{\small \texttt{theta\_cm0(i,"dd",r)}}].

The price index for subsistence demand is defined by a Leontief cost function:
\begin{equation}
p^{SD}_r = \sum_i \theta_i p^c_{ir},
\end{equation}
in which $\theta_i$ represents the commodity $i$ share of subsistence demand \hfill \break [{\small \texttt{theta\_cm0(i,"sd",r)}}].

Finally, \GAMS code includes a third representation of final demand in the form of the \CDE demand system \citep{Hanoch1975}.
This representation, details of which are presented in Appendix D of \citet{LanzRutherford2016}, has been part of the \GTAP model in \textsc{gempack} for many years, and our implementation follows those of \citet{HertelvanderMensbrugghe2015} and \citet{Chen}.
Similar to the LES demand system, this approach permits calibrating empirical evidence on both own-price and income elasticities.

\paragraph{\bf Government demand (\texttt{Y("g",r)})}

Public expenditure is a fixed-coefficient aggregate of Armington composite goods.
Within each composite domestic and imported goods trade off with a constant elasticity of substitution.
The unit price indices for domestic and imported goods are given by:
$$p^{dg}_{ir} = p^Y_{ir} \frac{1+t^{dg}_{ir}}{1+\overline{t}^{gd}_{ir}} $$
and
$$p^{ig}_{ir} = p^M_{ir} \frac{1+t^{ig}_{ir}}{1+\overline{t}^{ig}_{ir}}$$
The composite price of the $i$th good is then:
$$pg_{ir} = \left( \theta (p^{dg}_{ir})^{1-\sigma} + (1-\theta)
(p^{ig}_{ir})^{1-\sigma} \right)^{1/(1-\sigma)},$$ in which $\theta$ represents the domestic share of public demand [{\small \texttt{theta\_vdfm(i,"g",r)}}].

The unit cost of public services ($G_{r}$) is defined by the Leontief cost coefficients:
\begin{equation}
\sum_i \theta_i pg_{ir}  = p^G_r,
\end{equation}
in which $\theta_i$ represents the value share of commodity $i$ in public expenditure \hfill \break [{\small \texttt{theta\_cm0(i,"g",r)}}].

\paragraph{\bf Aggregate imports (\texttt{M(i,r)})}

An import cost index applies export taxes, trade and transport margins, and import tariffs to the producer supply prices in exporting regions:
$$py^m_{isr} = p^Y_{is}
\frac{(1-t^{xs}_{isr})(1+t^{ms}_{isr})} {(1-\overline{t}^{xs}_{isr})(1+\overline{t}^{ms}_{isr})}\,.$$ When $\textit{etrndx}_{i}<\infty$, the supply price $p^Y_{is}$ is replaced by export price $p^X_{is}$.
The unit price of transportation services is given by: $$pt^m_{jisr} = p^T_j \frac{1+t^{ms}_{isr}}{1+\overline{t}^{ms}_{isr}}\,.$$

Transportation margins enter as fixed coefficients with bilateral trade flows, so the unit delivered price is a convex combination of the unit prices with weights corresponding to base year value shares:
$$pyt^m_{isr} = \theta py^m_{isr} + \sum_j \theta^T_j pt^m_{jisr},$$
in which $\theta$ represents good $i$ share of imports [{\small \texttt{theta\_vxmd(i,s,r)}}]; and $\theta^T_j$ represents the value share of transportation service $j$ in the import price \hfill \break [{\small \texttt{theta\_vtwr(j,i,s,r)}}].

Having formed a price index for bilateral imports from region $s$ to region $r$, the \CES cost index can be defined on the basis of value shares and the elasticity of substitution across imports from different regions, $\sigma = \textit{esubm}_i$:
$$\textit{cim}_{ir} = \left(\sum_s \theta_s (pyt^m_{isr})^{1-\sigma)}
\right)^{1/(1-\sigma)} ,$$ in which $\theta_s$ is the value share of bilateral imports from region $s$ [{\small \texttt{theta\_m(i,s,r)}}].

The import activity ($m_{ir}$) has a zero profit condition which relates the unit cost of imports to the market price of the import aggregate:
\begin{equation}
cim_{ir} = p^M_{ir}\,.
\end{equation}

\paragraph{\bf International transportation services (\texttt{YT(j)})}

For simplicity, the unit cost of a transportation service depends on the benchmark value shares of region-specific services through a Cobb-Douglas cost function.
Under perfect competition with free entry, the unit cost of international transport services equals the equilibrium market price:
\begin{equation}
\prod_r (p^Y_{jr})^{\theta_r} = p^T_j
\end{equation}
in which $\theta_r$ represents the region $r$ share of transportation service $j$ [{\small \texttt{theta\_vst(j,r)}}].
When $\textit{etrndx}_{i}<\infty$, the supply price $p^Y_{is}$ in this equation is replaced by export price $p^X_{is}$.

\paragraph{\bf Sector-specific factor transformation (\texttt{FT(f,r)})}

The unit value of sector-specific factors is defined as a \CET aggregate of returns to factor $f$ across sectors $j$:
$$\textit{pvfm}_{fr} = \left( \sum_j \theta_j ps_{fjr}^{1+\eta}
\right)^{1/(1+\eta)} \quad f \in s_{f}\,.$$ in which $\theta_j$ is the sector $j$ share of earnings for factor $f$ [{\small \texttt{theta\_evom(f,j,r)}}].

The constant elasticity of transformation frontier defines the profit-maximizing allocation of factors to individual sectors.
In equilibrium, the unit value of the aggregate factor is equal to the maximum unit earnings:
\begin{equation}
p^S_{fr} = \textit{pvfm}_{fr} \quad f \in s_f
\end{equation}

\subsubsection{Market clearance}

Supply-demand conditions apply to all goods and factors.
Benchmark demand and supply quantities appear as scale factors in many of these equations, typically multiplied by activity levels which are equal to unity in the reference equilibrium.\footnote{While not crucial for representation of the model as a nonlinear system of equations, we follow the \MCP convention in writing out the market clearance conditions.
The equations are ``oriented'', with supply variables on the \LHS and demands on the \RHS.
Hence, the sense of the equation is $supply \geq demand$.
In the core model equilibrium prices should always be positive, but in extensions of the standard model it might be quite common to introduce inequalities and complementary slackness, in which case the proper orientation of the equations is essential.
Hence, in equilibrium should the price of a good be zero, economic equilibrium is then consistent with a market in which $supply > demand$.}

\paragraph{\bf Firm output (\texttt{P(i,r)})}

Aggregate output of good $i$ in region $r$ in the reference equilibrium is \texttt{vom(i,r)}:
\begin{equation}
\label{Youtput}
Y_{ir} \textit{vom}_{ir} =
\sum_j \textit{ddfm}_{ijr} + \textit{ddfm}_{iCr} + \textit{ddfm}_{iIr}
+ \textit{ddfm}_{iGr} + \sum_s \textit{dxmd}_{irs} + \textit{dst}_{ir}
\end{equation}
where the compensated demand functions can be obtained by differentiating the unit cost functions:
\begin{eqnarray}
\textit{ddfm}_{ijr} &=& Y_{jr} \textit{vdfm}_{ijr} \left( \frac{\textit{ci}_{ijr}}{p^d_{ijr}} \right)^\sigma \nonumber\\
\textit{ddfm}_{iCr}  &=& C_r \left( \textit{vdfm}_{iCr}\frac{p^C_r}{pc_{ir}} \right) \left(\frac{p^C_{ir}}{p^{\textit{dc}}_{ir}} \right)^\sigma \nonumber\\
\textit{ddfm}_{iIr}  &=& I_r \textit{vdfm}_{iIr} \nonumber\\
\textit{ddfm}_{iGr} &=& G_r \textit{vdfm}_{iGr}\left(\frac{pg_{ir}}{p^{\textit{dg}}_{ir}} \right)^\sigma   \nonumber\\
\textit{dxmd}_{isr} &=& M_{ir} \textit{vxmd}_{isr} \left(\frac{p^M_{ir}}{\textit{pyt}^m_{isr}} \right)^\sigma \nonumber\\
\textit{dst}_{jr} &=& \textit{YT}_j  \textit{vst}_{jr} \frac{p^T_j}{p^Y_{jr}}\nonumber
\end{eqnarray}

When $\textit{etrndx}_{i}<\infty$, equation (\ref{Youtput}) is replaced by two equations, one representing the market for domestic output and another representing the market for exports:
\begin{eqnarray}
Y_{ir} \textit{ddm}_{ir} &=&
\sum_j \textit{ddfm}_{ijr} + \textit{ddfm}_{iCr} + \textit{ddfm}_{iIr}
+ \textit{ddfm}_{iGr} \label{Doutput}\\
Y_{ir} \textit{dxm}_{ir} &=&
\sum_s \textit{dxmd}_{irs} + \textit{dst}_{ir}
\,.\label{Xoutput}
\end{eqnarray}
In these expressions, the compensated domestic and export supply coefficients are given by:
\begin{eqnarray}
\textit{ddm}_{ir} &=& \textit{vdm}_{ir} \left( \frac{ r^Y_{ir} } {
p^D_{ir}} \right)^{\eta} \nonumber \\
\textit{dxm}_{ir} &=& \textit{vxm}_{ir} \left( \frac{ r^Y_{ir} } {
p^X_{ir}} \right)^{\eta}. \nonumber
\end{eqnarray}

\paragraph{\bf Private consumption (\texttt{P("c",r)})}

Consumer demand in region $r$ in the reference equilibrium is $\textit{vom}_{Cr}$ hence:
\begin{equation}
C_r \textit{vom}_{Cr}  = \frac{\textit{RA}_r}{p^C_r}
\end{equation}

\paragraph{\bf Composite imports (\texttt{PM(i,r)})}

The aggregate value of imports of good $i$ in region $r$ in the reference equilibrium is $\textit{vim}_{ir}$:
\begin{equation}
M_{ir} \textit{vim}_{ir} = \sum_j \textit{difm}_{ijr} +
\textit{difm}_{iCr} + \textit{difm}_{iGr}
\end{equation}
where compensated demand functions are given by:
\begin{eqnarray}
\textit{difm}_{ijr} &=& Y_{jr} \textit{vifm}_{ijr} \left( \frac{ci_{ijr}}{p^i_{ijr}} \right)^\sigma\nonumber\\
\textit{difm}_{iCr} &=& C_r \textit{vifm}_{iCr}  \left(\frac{pc_{ir}}{p^{ic}_{ir}}\right)^\sigma  \frac{p^C_r}{\textit{pc}_{ir}} \nonumber\\
\textit{difm}_{iGr} &=& G_r \textit{vifm}_{iGr} \left(\frac{pg_{ir}}{p^{ig}_{ir}}\right)^\sigma \nonumber
\end{eqnarray}

\paragraph{\bf Transport services (\texttt{PT(j)})}

The aggregate demand (and supply) for transport service $j$ in the benchmark equilibrium is $\textit{vtw}_j$:
\begin{equation}
\textit{YT}_j  \textit{vtw}_j = \sum_{isr} \textit{dtwr}_{jisr}
\end{equation}
where
$$\textit{dtwr}_{jisr} = M_{ir} \textit{vtwr}_{jisr}  \left(
\frac{p^M_{ir}}{\textit{pyt}^m_{isr}} \right)^\sigma\,.$$

\paragraph{\bf Primary factors (\texttt{PF(f,r)})}

The aggregate demand (and supply) of primary factor $f$ in region $r$ is $\textit{evom}_{fr}$:
\begin{equation}
\textit{evom}_{fr} = \left\{
\begin{array}{lr}
\sum_j \textit{dfm}_{fjr} & f \in m_f \\
\textit{evom}_{fr} \textit{FT}_{fr} & f \in s_f
\end{array} \right.
\end{equation}
where the demand for primary factor is given by:
$$\textit{dfm}_{fjr}  = Y_{jr} \textit{vfm}_{fjr}
\left( \frac{c^f_{jr}}{p^{pf}_{fjr}}\right)^\sigma\,.$$

\paragraph{\bf Specific factors (\texttt{PS(f,j,r)})}

The net value of benchmark payments to factor $f$ in sector $j$ in region $r$ is \texttt{vfm(f,j,r)}:
\begin{equation}
\textit{vfm}_{fjr} \left( \frac{p^S_{s_{f},jr}}{p^F_{s_{f},r}} \right)^\eta =
\textit{dfm}_{fjr}\end{equation}
where the demand for primary factor is written above.

\subsubsection[Regional budget]{Regional budget (\texttt{RA(r)})}

Private and public incomes are given by :
\begin{equation}
\textit{RA}_r = \sum_f p^F_{fr} \textit{evom}_{fr} + p^I_n
\textit{vb}_r -  p^I_{r} \textit{vom}_{Ir} - p^G_{r}
\textit{vom}_{Gr}  + \mathcal{R}_r
\end{equation}
The base year current account deficit in region $r$ is \texttt{vb(r)}, and region $r=n$ corresponds to the ``numeraire region'' whose investment price (\texttt{P("i",rnum)} in the code) denominates international capital flows.
The numeraire region is the modelled region with the largest consumption (macro \texttt{loadrm}); in a single-region small open economy closure it is the modelled region itself, so that the current account is valued at its domestic investment price, the rest of the world providing the corresponding investment goods (following conventional static trade theory, we hold the current account deficit fixed in counterfactual analysis).
Furthermore, tax revenue in region $r$ consists of output taxes, intermediate demand taxes, factor taxes, final demand taxes, import tariffs, and export subsidies:
\begin{equation}
\begin{array}{ll}
\mathcal{R}_r = & \sum_j \mathcal{R}^o_{jr}
                 + \sum_{ij} \left( \mathcal{R}^{fd}_{ijr} +
                 \mathcal{R}^{fi}_{ijr} \right)
                 + \sum_{fj} \mathcal{R}^f_{fjr} \\
                &+ \sum_{i}  \left( \mathcal{R}^{pd}_{ir} +
                 \mathcal{R}^{pi}_{ir} + \mathcal{R}^{gd}_{ir} +
                 \mathcal{R}^{gi}_{ir}  \right)
                 - \sum_{is} \mathcal{R}^{xs}_{irs}
                 + \sum_{is} \mathcal{R}^{ms}_{isr}
\end{array}
\end{equation}
Each of these components of tax revenue can be calculated as an ad-valorem or proportional tax rate times a market price times the quantity demanded or produced.

Taxes related to $Y_{jr}$ include output taxes:\footnote{Tax revenues in the GAMS codes -- both \MCP and \MGE are represented by the \textit{macros} indicated in square brackets.}
$$\mathcal{R}^o_{jr} =  t^o_{jr} \textit{vom}_{jr} p^Y_{jr}
Y_{jr}\, \quad [\footnotesize{\texttt{REV\_TO(g,r)}}]\,,$$ tax revenue from intermediate inputs:
$$\mathcal{R}^{fd}_{ijr} =  t^{fd}_{ijr} p^Y_{ir}
\textit{ddfm}_{ijr}\, \quad [\footnotesize{\texttt{REV\_TFD(i,j,r)}}]\,,$$
$$\mathcal{R}^{fi}_{ijr} = t^{fi}_{ijr} p^M_{ir}
\textit{difm}_{ijr}, \quad [\footnotesize{\texttt{REV\_TFI(i,j,r)}}]\,,$$ and factor tax revenue:
$$\mathcal{R}^f_{fjr} = t^{f}_{fjr} p^F_{fr} \textit{dfm}_{fjr} \quad
[\footnotesize{\texttt{REV\_TF(f,g,r)}}]\,.$$

Taxes on household consumption of domestic and imported goods are:
$$\mathcal{R}^{pd}_{ir} = t^{pd}_{ir} p^Y_{ir} \textit{ddfm}_{iCr},
\quad [\footnotesize{\texttt{REV\_TFD(i,"C",r)}}]\,,$$ and
$$\mathcal{R}^{pi}_{ir} = t^{pi}_{ir} p^M_{ir} \textit{difm}_{iGr}
\quad [\footnotesize{\texttt{REV\_TFI(i,"C",r)}}]\,.$$

Taxes on public demand for domestic and imported goods are:
$$\mathcal{R}^{gd}_{ir} = t^{gd}_{ir} p^Y_{ir} \textit{ddfm}_{iGr}
\quad [\footnotesize{\texttt{REV\_TFD(i,"G",R)}}]\,,$$ and
$$\mathcal{R}^{gi}_{ir} =  t^{gi}_{ir} p^M_{ir} \textit{difm}_{iGr}
\quad [\footnotesize{\texttt{REV\_TFI(i,"G",r)}}]\,.$$

Export subsidies (paid by the government in the exporting region) are:
$$\mathcal{R}^{xs}_{irs} = t^{xs}_{irs} p^Y_{ir}
\textit{dxmd}_{irs} \quad [\footnotesize{\texttt{REV\_TXS(i,r,s)}}]\,,$$ and import tariff revenues are given by:
$$\mathcal{R}^{ms}_{isr} = t^{ms}_{isr} \left(p^Y_{is}
(1-t^{xs}_{isr}) \textit{dxmd}_{isr} + \sum_j p^T_j \textit{dtwr}_{jisr} \right) \quad [\footnotesize{\texttt{REV\_TMS(i,s,r)}}]\,.$$

\subsection[Closures]{Closure with respect to the rest of the world: \GMR, \SOE, and \SATE}\label{app:closures}

This subsection is new relative to \citet{LanzRutherford2016}.
It documents how the macro \texttt{loadrm(rr)} (file \texttt{code/loadmdl.gms}) and the flags \texttt{sate} and \texttt{sate2} (file \texttt{code/gtap12data.gms}) select the trade closure, and states the equilibrium conditions which are added to or replace those of Appendix \ref{equil}.
Let $R^M$ (\texttt{rm(r)}) denote the regions that are part of the model and $R^X$ (\texttt{rx(r)}) the remaining regions which form the rest of the world (\textsc{row}).
The \GMR model is obtained with $R^M$ equal to all regions and $R^X=\emptyset$; all conditions of Appendix \ref{equil} then apply to all regions and no additional variables are active.

\subsubsection{Variables and parameters}

Table \ref{tab:closurevars} lists the variables that are active only in the small open economy closures.
In addition, the following benchmark parameters are computed by \texttt{loadrm} for $r\in R^M$: exports to the rest of the world, $\textit{vem}_{ir}=\sum_{s\in R^X}\textit{vxmd}_{irs}$ [\texttt{vem(i,r)}]; the trade-weighted average export subsidy rate on these exports, $t^{xs,row}_{ir}=\sum_{s\in R^X}\textit{vxmd}_{irs}\,t^{xs}_{irs}/\textit{vem}_{ir}$ [\texttt{rtxs\_row(i,r)}], and the associated reference price $\textit{pem}^0_{ir}=1-\overline{t}^{xs,row}_{ir}$ [\texttt{pem0(i,r)}].
For $s\in R^X$, $\textit{vem}_{is}=\sum_{r\in R^M}\textit{vxmd}_{isr}$ denotes the exports of \textsc{row} region $s$ to the modelled regions.

\begin{table}[t!]
\begin{center}
\caption{Variables of the small open economy closures.}\label{tab:closurevars}
\small
\begin{tabular}{lp{6.2cm}lr}
\toprule
Variable & Definition & \GAMS variable & Active in \\
\midrule
$X_{ir}$, $r\in R^M$ & Exports of model region $r$ to \textsc{row} & \texttt{X(i,r)} & \SOE, \SATE \\
$X_{is}$, $s\in R^X$ & Exports of \textsc{row} region $s$ to the model & \texttt{X(i,s)} & \SOE, \SATE \\
$p^{FX}$ & Price of foreign exchange (real exchange rate) & \texttt{PFX} & \SOE, \SATE \\
$\textit{ROW}$ & Income of the rest-of-world agent & \texttt{ROW} & \SOE, \SATE \\
\addlinespace
$p^{PX}_{ir}$ & Price of exports of $r$ in its export market & \texttt{PPX(i,r)} & \texttt{sate} \\
$E_{ir}$ & Export demand index (benchmark $=1$) & \texttt{EXPORTS(i,r)} & \texttt{sate} \\
$\textit{RX}_{ir}$ & Foreign exchange earned from exports (index) & \texttt{REVEXP(i,r)} & \texttt{sate} \\
\addlinespace
$p^{FIX}_{ir}$ & Price of the \textsc{row}-owned fixed factor in export demand & \texttt{PFIX(i,r)} & \texttt{sate2} \\
\bottomrule
\end{tabular}
\end{center}
\end{table}

\subsubsection{Small open economy with fixed world prices (\SOE)}

Exports of a modelled region to the rest of the world are converted into foreign exchange at a fixed rate.
With $p^X_{ir}$ denoting the export supply price (equal to $p^Y_{ir}$ if $\textit{etrndx}_i=\infty$), the zero profit condition of $X_{ir}$ reads
\begin{equation}\label{zpX}
p^X_{ir}\,\frac{1-t^{xs,row}_{ir}}{1-\overline{t}^{xs,row}_{ir}} \;\geq\; p^{FX}
\qquad \perp \; X_{ir}\geq 0, \quad r\in R^M,
\end{equation}
so that the world market absorbs any quantity at a price which is fixed in terms of foreign exchange.
Symmetrically, \textsc{row} supplies the varieties it exports to the modelled regions at a fixed price, $p^E_{is}=p^{FX}$ for $s\in R^X$, and these varieties enter the Armington import aggregates of the modelled regions as in Appendix \ref{equil}.
The market for foreign exchange clears through $p^{FX}$:
\begin{equation}\label{mktFX}
\sum_{r\in R^M}\textit{pem}^0_{ir}\textit{vem}_{ir}X_{ir} + \textit{rowpfx}
= \sum_{s\in R^X}\textit{vem}_{is}X_{is} + \sum_{j,s\in R^X}\textit{vst}_{js}\,\textit{YT}_j
+ \frac{\textit{ROW}}{p^{FX}},
\end{equation}
where the rest-of-world agent is endowed with an amount of foreign exchange, \textit{rowpfx} [\texttt{rowpfx}], provides the transport services that \textsc{row} regions export, and finances the current account balances of the modelled regions.
Its income is $\textit{ROW}=p^{FX}\textit{rowpfx} - p^I_n\sum_{r\in R^M}\textit{vb}_r - \sum_j p^T_j\sum_{i,s}\sum_{k\in R^X}\textit{vtwr}_{jisk}$ (the last term being the transport services used in shipments to \textsc{row} regions), which it spends on foreign exchange.
Equations \eqref{zpX}--\eqref{mktFX} are the closure of \citet{LanzRutherford2016}.
Without a finite elasticity of transformation between domestic and export supply ($\textit{etrndx}_i=\infty$, the \texttt{gtap} model version), \eqref{zpX} pins down the domestic producer price of every exported good, and output responds to trade policy in a bang-bang fashion.

\subsubsection{Small Armington trade economy, implementation \texttt{sate}}

In the \SATE closure, the export activity sells into a separate export market with price $p^{PX}_{ir}$,
\begin{equation}\label{zpXsate}
p^X_{ir}\,\frac{1-t^{xs,row}_{ir}}{1-\overline{t}^{xs,row}_{ir}} \;\geq\; p^{PX}_{ir}
\qquad \perp \; X_{ir}\geq 0,
\end{equation}
where the world's demand for the region's exports has the constant price elasticity $\sigma^m_i=\textit{esubm}_i$:
\begin{equation}\label{exportdemand}
E_{ir} = \left(\frac{p^{FX}}{p^{PX}_{ir}}\right)^{\sigma^m_i},
\qquad\quad
\textit{pem}^0_{ir}\textit{vem}_{ir}X_{ir} = \textit{pem}^0_{ir}\textit{vem}_{ir}E_{ir}
\quad \perp\; p^{PX}_{ir}.
\end{equation}
The representative agent of region $r$ sells the exports and receives their value in foreign exchange,
\begin{equation}\label{revexp}
\textit{RX}_{ir}\,p^{FX} = E_{ir}\,p^{PX}_{ir},
\end{equation}
so that the term $\sum_i \textit{pem}^0_{ir}\textit{vem}_{ir}\left(p^{FX}\textit{RX}_{ir}-p^{PX}_{ir}E_{ir}\right)=0$ is added to $\textit{RA}_r$ and the supply of foreign exchange in \eqref{mktFX} becomes $\sum_{r\in R^M}\textit{pem}^0_{ir}\textit{vem}_{ir}\textit{RX}_{ir}$.
In \MPSGE, $E_{ir}$ and $\textit{RX}_{ir}$ are auxiliary variables which ration an endowment of $-\textit{pem}^0_{ir}\textit{vem}_{ir}$ units of export goods and $+\textit{pem}^0_{ir}\textit{vem}_{ir}$ units of foreign exchange of agent \texttt{RA(r)} (see the \texttt{\$demand:RA(r)} blocks and the \texttt{\$constraint:EXPORTS} and \texttt{\$constraint:REVEXP} blocks in the code listing below).
Equations \eqref{zpXsate}--\eqref{revexp} require no information about the accounts of the rest of the world: the export demand curve is fully characterized by benchmark exports and the Armington elasticity.
Imports remain available at fixed world prices.

\subsubsection{Small Armington trade economy, implementation \texttt{sate2}}

Implementation \texttt{sate2} follows the large open economy formulation of \citet{Schreiber2024}.
The export activity combines the region's exports with a fixed factor owned by the rest of the world in a Cobb-Douglas technology and produces foreign exchange:
\begin{equation}\label{zpXsate2}
\left(p^X_{ir}\,\frac{1-t^{xs,row}_{ir}}{1-\overline{t}^{xs,row}_{ir}}\right)^{1-\theta^{x}_{i}}
\left(p^{FIX}_{ir}\right)^{\theta^{x}_{i}} \;\geq\; p^{FX}
\qquad \perp \; X_{ir}\geq 0,
\end{equation}
with the cost share of the fixed factor set to $\theta^x_i=1/\sigma^m_i$ [\texttt{theta\_xi(i)}] and a benchmark endowment of the fixed factor of $\textit{fix}^0_{ir}=\theta^x_i/(1-\theta^x_i)\,\textit{pem}^0_{ir}\textit{vem}_{ir}$ [\texttt{fix0(i,r)}], which accrues to the \textsc{row} agent.
With $p^{FX}$ given, the derived demand for the region's exports has price elasticity $1/\theta^x_i=\sigma^m_i$ \citep[Proposition 2]{Schreiber2024}, so that \texttt{sate} and \texttt{sate2} describe the same export demand curve.
The two implementations differ only in bookkeeping: in \texttt{sate2}, the rents associated with the downward sloping export demand curve are explicitly attributed to the rest of the world.
In the simulations reported in this paper, both implementations give identical results (up to the six digits reported).

\subsubsection{Average export taxes on exports to the rest of the world}

In the code of \citet{LanzRutherford2016}, $t^{xs,row}_{ir}$ is computed once by \texttt{loadrm} and does not follow subsequent changes of the bilateral rates $t^{xs}_{irs}$.
Program \texttt{tgrid.gms} now provides the macro \texttt{updrow(rr)} which recomputes $t^{xs,row}_{ir}$ after each change of export taxes (environment variable \texttt{--soextax=yes}, the default).
As a consequence, in the small open economy closures the counterfactual scaling of export taxes affects exports to the rest of the world, as it does in the \GMR model.
In the 2016 paper, only import tariffs moved in the \SOE simulations.

\subsection{Calibration of trade elasticities}\label{app:thetadx}

Our \SOE model is based on the idea that goods are differentiated both on the basis of origin and destination.
In this setting, production supplies goods for the domestic and export markets with a technology capturing the idea that the quantities supplied to these submarkets depend on relative prices.
Each traded good then has three associated trade elasticities: i) $\texttt{etadx(i)}$, the elasticity of transformation between goods produced for the the domestic and export markets; ii) $\texttt{esubdm(i)}$, the elasticity of substitution between domestic and imported varieties entering intermediate and final demand, and iii) $\texttt{esubm(i)}$, the elasticity of substitution between imports from different countries.

In the \GMR model corresponding to the standard \GTAP model, $\eta^{DX} = \infty$, so goods produced for the domestic and export markets are perfect substitutes.
In the \GTAPINGAMS framework, $\eta^{DX}$ can take on any nonnegative value.
A finite value is typically specified in the \SOE model.
When doing this, we adjust the domestic-import demand elasticity to maintain \textit{consistency} with trade responses in the standard \GTAP model.
We interpret consistency from a Marshallian perspective.
In a given market, suppose the import demand elasticity is $\epsilon$,\footnote{When value shares are small, the price elasticity of demand for imports is approximately equal to $-\texttt{esubdm(i)}$.} and benchmark rate of protection is $\bar t$.
Scaling units so that the FOB supply price is $1$, the demand price is $1+\bar t$, and the quantity demanded in the benchmark is $1$.

When the import supply elasticity from partner countries, $\eta$, is infinite, the import supply price is fixed ($p=1$).
Therefore, if the tariff was removed and demand is isoelastic, the new equilibrium quantity and price (in an isoelastic model) are given by:
\begin{equation}
\label{newqty}
q^* = \left(1 + \bar t \right)^{|\epsilon|}\,.
\end{equation}
With a fixed import supply price, the quantity demanded at $p=1$ is the new equilibrium.
This outcome is illustrated in Figure \ref{etainfty}, in which benchmark import demand is at price-quantity point $a$, benchmark import supply is at point $b$, and there is a benchmark trade barrier equal to 20\%.
The free-trade equilibrium is at point $c$, reflecting benchmark import demand elasticity $\epsilon = 5$.

\begin{figure}[p!]

\begin{subfigure}[b]{0.9\textwidth}
\centering
\includegraphics[height=0.4\textheight]{jgea_new/figures/figure9a.tikz}
\caption{ $\epsilon=5$, $\eta=\infty$.}
\label{etainfty}
\end{subfigure}

\vspace{0.2in}

\begin{subfigure}[b]{0.9\textwidth}
\centering
\includegraphics[height=0.4\textheight]{jgea_new/figures/figure9b.tikz}
\caption{$\epsilon=16.7$, $\eta=7.1$.}
\label{theta30}
\end{subfigure}

\caption{Consistent Calibration of $\epsilon$ and $\eta$.}

\end{figure}

Now instead of a perfectly elastic import supply we consider $\eta < \infty$.
In this case, removing tariffs implies a new equilibrium supply price $p^*$, and the associated equilibrium supply and demand $q^*$ from (\ref{newqty}) is given by:
$$q^* = \left( \frac{1+\bar t}{p^*} \right)^{|\epsilon|} =
\left( p^* \right)^\eta\,.$$ The aggregate trade response under $\eta < \infty$ equals the response with $\eta = \infty$ when the demand elasticity ($\epsilon$) is adjusted upward:
$$p^* = (1 + \bar t)^{|\epsilon|/\eta}\,,$$
the reparametrized elasticity of demand is given by:
$$|\epsilon|  = \frac{|\epsilon|}{1 - |\epsilon|/\eta}\,.$$

Let $\theta = 1 - |\epsilon|/\eta$.
Consistency of trade quantity responses in the revised model implies that $\theta \geq 0$.
When $\eta = \infty$, then we have the original model with a horizontal import supply curve.
Setting $\theta=0.3$, the new Marshallian model appears in Figure \ref{theta30}.
As in the previous model, $a$ and $b$ are benchmark demand and supply with a 20\% tariff barrier.
While free trade in the original model is at point $c$, free trade in the revised model is at point $d$.
Supply is less elastic, so social surplus increases are captured by the supply side of the market, and the equilibrium price following tariff removal increases by over 10\% from the benchmark level.

Note that, with this calibration, \texttt{thetadx}$=\theta$ applies in both directions of trade and to all closures: the elasticity of transformation between domestic and export supply is $\eta^{DX}_i=\sigma^D_i/(1-\theta)$ and the elasticity of substitution between domestic and imported goods is raised to $\sigma^{DM}_i=\sigma^D_i/\theta$, where $\sigma^D_i$ is the \GTAP value \texttt{esubd(i)}.
For $\theta=1$ the standard \GTAP specification is obtained ($\eta^{DX}_i=\infty$, $\sigma^{DM}_i=\sigma^D_i$); this is the \texttt{gtap} model version of the present paper.
For $\theta=0.5$, the default of \citet{LanzRutherford2016} and the \texttt{gtapingams} model version of this paper, $\eta^{DX}_i=\sigma^{DM}_i=2\sigma^D_i$.

In the \GAMS code, this adjustment takes place near the end of \texttt{gtap12data.gms}; the default \texttt{thetadx}$=0.5$ is set in \texttt{mge.gms} and may be overridden on the command line.
The run-time environment variable \texttt{thetadx} defines how the adjustment of export supply and import demand is to be calibrated to match the \GMR model.
The associated assignments are reported here.

\begin{Verbatim}
parameter       etadx(g)        Export elasticity of transformation D vs E,
                esubdm(i)       Elasticity of substitution D vs M,
                vxm(g,r)        Export market supply;

$if not set thetadx $set thetadx 0.5

etadx(i)$(1-%thetadx%) = esubd(i)/(1-%thetadx%);
etadx(i)$(not (1-%thetadx%)) = +inf;
esubdm(i) = esubd(i)/%thetadx%;
vxm(i,s) = (vst(i,s) + sum(r,vxmd(i,s,r)))$(1/etadx(i));
\end{Verbatim}

%\clearpage
\subsection[Programs]{Programs and workflow of \GTAPtwelve}\label{app:programs}

The distribution keeps the folder structure of \citet[Appendix E]{LanzRutherford2016}: \texttt{code/} contains the model and data routines, \texttt{defines/} the aggregation maps, \texttt{build/} the working directory for dataset construction and model checks, \texttt{gamsdata/20yy/} the aggregated datasets, and \texttt{forensics/} the programs for the illustrative calculations.
Instead of the \GAMS command scripts \texttt{build.gms} and \texttt{run.gms}, the sequence of programs is managed by the Python tool \texttt{doit} (file \texttt{dodo.py}), which re-runs a step only if its inputs have changed.
Table \ref{tab:dodo} lists the tasks.

\begin{table}[ht!]
\begin{center}
\caption{Tasks of the \texttt{doit} build stream (file \texttt{dodo.py}).}\label{tab:dodo}
\small
\begin{tabular}{lp{9.5cm}}
\toprule
Task & Action \\
\midrule
\texttt{readgdx} & Reads the \GDX version of the \GTAP 12 Data Base
(\texttt{GDX12a\_AY16994.zip}, all reference years) and writes
\texttt{gsd.gdx} for the chosen reference year (replaces
\texttt{flex2gdx.gms}, which translated \HAR files). \\
\texttt{filter} & Filters small values (\texttt{--nd=5}) and
rebalances the dataset (\texttt{filter.gms}). \\
\texttt{build} & Aggregates to \texttt{gtap12ingams}, \texttt{g20},
\texttt{g20\_iea} and \texttt{g20\_macro} (\texttt{gtapaggr.gms}) and
recalibrates \CDE demand (\texttt{cdecalib.gms}). \\
\texttt{test} & Benchmark replication, \MCP/\MGE consistency, and \CDE
elasticity checks (\texttt{bmkchk.gms}, \texttt{mcpmge.gms},
\texttt{cdechk.gms}). \\
\texttt{grid} & Runs all counterfactuals of Section \ref{sec:comparison}
(\texttt{forensics/ssasolve.gms} calling \texttt{tgrid.gms}). \\
\texttt{merge} & Merges the results (\texttt{ssamerge.gms}) into
\texttt{results.gdx} and an Excel pivot table. \\
\texttt{tarifffigs} & Benchmark trade taxes by commodity and partner
(\texttt{tradetax.gms}). \\
\texttt{resultsfigs} & Writes the data files for the figures
(\texttt{process\_results.gms}). \\
\bottomrule
\end{tabular}
\end{center}
\end{table}

Program \texttt{forensics/tgrid.gms} performs the sequence of counterfactuals for one region.
Its environment variables are listed in Table \ref{tab:tgrid}; for example,
\begin{Verbatim}
gams tgrid --ds=g20_macro --rcalc=ZAF --mdl=sate --thetadx=1
\end{Verbatim}
computes the \SATE results of the \texttt{gtap} model version for South Africa.

\begin{table}[ht!]
\begin{center}
\caption{Environment variables of \texttt{forensics/tgrid.gms}.}\label{tab:tgrid}
\small
\begin{tabular}{lp{8cm}l}
\toprule
Variable & Meaning & Default \\
\midrule
\texttt{--ds} & Dataset in \texttt{gamsdata/20yy/} & \texttt{g20\_macro} \\
\texttt{--yr} & Reference year & \texttt{23} \\
\texttt{--rcalc} & Region whose trade taxes are scaled & \texttt{CHN} \\
\texttt{--mdl} & Trade closure: \texttt{gmr}, \texttt{soe}, \texttt{sate}, \texttt{sate2} & \texttt{sate2} \\
\texttt{--thetadx} & $1$: \texttt{gtap} version ($\eta^{DX}=\infty$); $0.5$: \texttt{gtapingams} version & \texttt{1} \\
\texttt{--fsubm} & Scaling factor for $\sigma^m_i$ (\texttt{esubm}); $0.5$ for \GMRLS & \texttt{1} \\
\texttt{--soextax} & Update $t^{xs,row}_{ir}$ when export taxes change (macro \texttt{updrow}) & \texttt{yes} \\
\bottomrule
\end{tabular}
\end{center}
\end{table}

The model is solved with the \MPSGE formulation (\texttt{code/mge.gms}); the \SATE closures have not been added to the algebraic \MCP formulation (\texttt{code/mcp.gms}), which remains available for the \GMR and \SOE closures.


\subsection[Code]{\MPSGE code of the model (\texttt{code/mge.gms})}\label{app:code}

The listing below is read directly from the file \texttt{code/mge.gms} of the distribution when this document is compiled.

\VerbatimInput[fontsize=\tiny,frame=single]{../code/mge.gms}

\section{Benchmark Trade Taxes}\label{app:tariffs}

Figures \ref{fig:taxcom1} to \ref{fig:taxpar2} summarize the commodity and partner structure of the taxes on imports into the four countries in the \texttt{g20} aggregation, in the style of \citet{LanzRutherford2016}.
They show import tariff rates of the four importing countries and export tax rates that the exporting regions charge on exports to the four countries, together with the tariff revenue of the four countries (in billion USD).
Averaged over all imports, the United States levy import tariffs of 1.1\% and face export taxes of their trading partners of 0.5\%.
For China, the corresponding figures are 1.9\% and 0.5\%, for South Africa 4.1\% and 1.4\%, and for Argentina 4.9\% and 1.3\%.
Protection is far from uniform.
In the United States, the highest rates apply to sugar, leather, dairy products, and wearing apparel; in China, to wheat, sugar, other cereals, and wool.
South Africa and Argentina levy rates of 15--32\% on wearing apparel, leather, and textiles.
Across partners, preferential agreements are visible: the United States levy almost no tariffs on imports from Mexico and Canada, Argentina none on imports from Brazil.
Tariff revenue amounts to 43.7 billion USD in the United States, 53.8 billion USD in China, 5.4 billion USD in South Africa, and 4.3 billion USD in Argentina.
It is concentrated on a few goods and partners.
Wearing apparel and leather account for 38\% of the tariff revenue of the United States, and motor vehicles for 14\% of the tariff revenue of China and 33\% of that of South Africa.
Imports from China account for 34--43\% of the tariff revenue of the United States, South Africa, and Argentina.

\begin{figure}[p!]
\centering
\input{jgea_new/figures/tariff_commodity_USA}
\vspace{0.1in}
\input{jgea_new/figures/tariff_commodity_CHN}
\pgfplotslegendfromname{tariffcom_legend}
\caption{Trade taxes on imports by commodity, United States and China (\GTAP 12, 2023; 23 most highly taxed commodities).}\label{fig:taxcom1}
\end{figure}

\begin{figure}[p!]
\centering
\input{jgea_new/figures/tariff_commodity_ZAF}
\vspace{0.1in}
\input{jgea_new/figures/tariff_commodity_ARG}
\pgfplotslegendfromname{tariffcom_legend}
\caption{Trade taxes on imports by commodity, South Africa and Argentina (\GTAP 12, 2023; 23 most highly taxed commodities).}\label{fig:taxcom2}
\end{figure}

\begin{figure}[p!]
\centering
\input{jgea_new/figures/tariff_txtm_USA}
\vspace{0.1in}
\input{jgea_new/figures/tariff_txtm_CHN}
\pgfplotslegendfromname{tariffpar_legend}
\caption{Trade taxes on imports by trading partner, United States and China (\GTAP 12, 2023).}\label{fig:taxpar1}
\end{figure}

\begin{figure}[p!]
\centering
\input{jgea_new/figures/tariff_txtm_ZAF}
\vspace{0.1in}
\input{jgea_new/figures/tariff_txtm_ARG}
\pgfplotslegendfromname{tariffpar_legend}
\caption{Trade taxes on imports by trading partner, South Africa and Argentina (\GTAP 12, 2023).}\label{fig:taxpar2}
\end{figure}

\end{appendices}

\newpage

\bibliographystyle{jgea}
\bibliography{jgea_new/jgea_new}

\end{document}
